% Options for packages loaded elsewhere
\PassOptionsToPackage{unicode}{hyperref}
\PassOptionsToPackage{hyphens}{url}
\PassOptionsToPackage{dvipsnames,svgnames,x11names}{xcolor}
\documentclass[
  british,
  10pt,
]{article}
\usepackage{xcolor}
\usepackage[margin=25mm]{geometry}
\usepackage{amsmath,amssymb}
\setcounter{secnumdepth}{5}
\usepackage{iftex}
\ifPDFTeX
  \usepackage[T1]{fontenc}
  \usepackage[utf8]{inputenc}
  \usepackage{textcomp} % provide euro and other symbols
\else % if luatex or xetex
  \usepackage{unicode-math} % this also loads fontspec
  \defaultfontfeatures{Scale=MatchLowercase}
  \defaultfontfeatures[\rmfamily]{Ligatures=TeX,Scale=1}
\fi
\usepackage{lmodern}
\ifPDFTeX\else
  % xetex/luatex font selection
\fi
% Use upquote if available, for straight quotes in verbatim environments
\IfFileExists{upquote.sty}{\usepackage{upquote}}{}
\IfFileExists{microtype.sty}{% use microtype if available
  \usepackage[]{microtype}
  \UseMicrotypeSet[protrusion]{basicmath} % disable protrusion for tt fonts
}{}
\makeatletter
\@ifundefined{KOMAClassName}{% if non-KOMA class
  \IfFileExists{parskip.sty}{%
    \usepackage{parskip}
  }{% else
    \setlength{\parindent}{0pt}
    \setlength{\parskip}{6pt plus 2pt minus 1pt}}
}{% if KOMA class
  \KOMAoptions{parskip=half}}
\makeatother
\usepackage{longtable,booktabs,array}
\newcounter{none} % for unnumbered tables
\usepackage{calc} % for calculating minipage widths
% Correct order of tables after \paragraph or \subparagraph
\usepackage{etoolbox}
\makeatletter
\patchcmd\longtable{\par}{\if@noskipsec\mbox{}\fi\par}{}{}
\makeatother
% Allow footnotes in longtable head/foot
\IfFileExists{footnotehyper.sty}{\usepackage{footnotehyper}}{\usepackage{footnote}}
\makesavenoteenv{longtable}
\ifLuaTeX
\usepackage[bidi=basic,shorthands=off]{babel}
\else
\usepackage[bidi=default,shorthands=off]{babel}
\fi
\ifLuaTeX
  \usepackage{selnolig} % disable illegal ligatures
\fi
\setlength{\emergencystretch}{3em} % prevent overfull lines
\providecommand{\tightlist}{%
  \setlength{\itemsep}{0pt}\setlength{\parskip}{0pt}}
\usepackage{microtype}
\usepackage{xurl}
\usepackage{etoolbox}
\emergencystretch=3em
\setlength{\parskip}{0.35em}
\AtBeginEnvironment{longtable}{\small}
\AtBeginEnvironment{verbatim}{\small}
\AtBeginDocument{\hypersetup{pdfsubject={Unpublished partial results and proof audits; neither target conjecture is proved}}}
\usepackage{bookmark}
\IfFileExists{xurl.sty}{\usepackage{xurl}}{} % add URL line breaks if available
\urlstyle{same}
\hypersetup{
  pdftitle={Local obstructions and proof audits for Casas--Alvero},
  pdflang={en-GB},
  colorlinks=true,
  linkcolor={blue},
  filecolor={Maroon},
  citecolor={Blue},
  urlcolor={Blue},
  pdfcreator={LaTeX via pandoc}}

\title{Local obstructions and proof audits for Casas--Alvero}
\usepackage{etoolbox}
\makeatletter
\providecommand{\subtitle}[1]{% add subtitle to \maketitle
  \apptocmd{\@title}{\par {\large #1 \par}}{}{}
}
\makeatother
\subtitle{Consolidated research report for review}
\author{}
\date{23 September 2026}

\begin{document}
\maketitle

{
\setcounter{tocdepth}{3}
\tableofcontents
}
\section{Scope and principal
conclusions}\label{scope-and-principal-conclusions}

This report consolidates an audit and an attempted solution of the
Casas--Alvero conjecture, with particular attention to unrestricted
degree 20. It supplies the resulting partial theorems, explicit
countermodels to proposed arguments, and finite certificates for review.
\textbf{Neither unrestricted degree 20 nor the conjecture in all degrees
is proved by this work.} No counterexample to the characteristic-zero
conjecture has been found.

For a field of characteristic zero, the conjecture says that a monic
polynomial \(f\) of degree \(n\geq2\) satisfying
\[\deg\gcd(f,f^{(i)})>0\qquad(1\leq i<n)\] must be \((X-a)^n\). The
common root may depend on \(i\). Our original two targets were this
statement for \(n=20\) without a sparsity restriction, and the statement
for every \(n\). Neither a proof gap in someone else's argument nor a
sparse special case fulfils either target.

The report has three mathematical parts. The first reconstructs the
proof audit and the seven-term lower bound in centered degree 20. The
second gives the subsequent local reductions and case exclusions for
general degree-20 coefficients. The third develops the valid
conclusions, and the explicit limitations, of several proposed
approaches in arbitrary degree. The supplied proof notes retain details
too long for the main narrative, including complete case lists,
polynomial identities, and separate internal audits.

\subsection{What has been established within the reviewed
arguments}\label{what-has-been-established-within-the-reviewed-arguments}

\begin{enumerate}
\def\labelenumi{\arabic{enumi}.}
\tightlist
\item
  The stated generality of Ghosh's Proposition 3.3 fails in
  characteristics five and seven. The explicit ideals, generator counts,
  and dependency analysis identify proof obligations; they do not
  disprove the characteristic-zero conjecture.
\item
  A hypothetical nontrivial degree-20 counterexample has at least seven
  nonzero terms after centering, counting its leading term. This is a
  statement about monomials, not distinct roots or derivative witnesses.
\item
  The characteristic-17 visible model has nine normalized nontrivial
  seeds over the full algebraic closure. One whole seed cannot lift.
  Eight remain. Three of the surviving branches have complete unramified
  marked lifting descriptions; a fourth has a complete description
  allowing ramification index at most two.
\item
  Complete coefficient-support systems in the row-9 branch and one in
  row 4 have been excluded by finite residue searches followed by exact
  lifting. The final coverage is 79 of 240 row-9 systems; 161 remain
  unexecuted. This is not a whole-branch exclusion.
\item
  In the row-8 branch, a valuation bound and a divided identity
  constrain all supports. A large first-scale stratum admits an exact
  finite classification of quartic leading models, an exact
  cluster-collapse conclusion, and a further conditional exclusion. The
  whole branch remains unresolved.
\item
  Finite formal and resultant-algebra descriptions provide exact
  remaining exclusion targets. Their unevaluated generic fibres are not
  empty merely because the rings are finite.
\item
  In arbitrary degree, valid moment and height bounds coexist with
  explicit countermodels to relaxed descent arguments. The repaired
  contraction argument still requires the missing Casas--Alvero step. An
  exact matrix representation is available, but its proposed eigenvector
  shortcut fails.
\end{enumerate}

These are internal mathematical conclusions supported by the arguments
and checks described below. They have not undergone external refereeing
or formal proof-assistant verification. Priority and publication
significance are not established. In particular, the earliest five-term
lower bound was found to follow from prior work; the report preserves
that correction rather than treating every intermediate bound as an
original contribution.

\subsection{Notation and interpretation of
computations}\label{notation-and-interpretation-of-computations}

Write \[f(X)=\sum_{j=0}^{n}\binom nj a_jX^{n-j},\qquad
G_j(X)=\sum_{i=0}^{j}\binom ji a_iX^{j-i}.\] Thus \(G_j\) is the monic
divided derivative of degree \(j\). Hasse order \(k\) means
\(H_kf=f^{(k)}/k!\) in characteristic zero, continued by the integral
binomial formula in positive characteristic. Deficiency index \(j\),
normalized derivative degree \(j\), and Hasse order \(n-j\) are related
but must not be interchanged. A support is the set of indices of exactly
nonzero normalized coefficients. An active coefficient can have zero
residue.

Local valuations are normalized by \(v(p)=1\). Ramified candidates can
have fractional positive valuation. Reduction in the residue field and
divisibility by \(p\) in the valuation ring are different assertions.
Whenever an identity is divided by \(p\), its required integral
divisibility is proved before division. Likewise, searches over a
specified finite extension are complete only after that extension has
been shown to contain every required algebraic residue root.

The computational assertions have three distinct levels. An enumeration
certifies a finite set of residue configurations. A finite-precision
lift certifies polynomial congruences. A separate coverage and
uniqueness argument explains why those congruences exclude
characteristic-zero candidates, including candidates originally lying in
ramified fields. The report does not infer the third level from
successful execution of the first two.

\subsection{Evidence organization}\label{evidence-organization}

\texttt{evidence/audit/} contains the original proof audit and its
certificates; \path{evidence/AUDIT_NOVELTY_CORRECTION.md} supersedes its
earliest novelty assessment. \texttt{evidence/seven-terms/} contains the
complete seven-term proof and its unchanged dependency packages.
\texttt{evidence/full/} contains the later unrestricted campaign,
including successful results and failed experiments. Historical files
can have outdated status language; this report and
\texttt{REVIEW\_GUIDE.md} identify what is superseded. The older files
are retained for traceability.

The top-level \texttt{PROOF\_INDEX.md} lists the proof notes, while
\texttt{replay.py} gives bounded, reproducible checks in a temporary
copy. The default replay does not launch exploratory solvers or
regenerate optional large discovery certificates. It records failures
and timeouts explicitly. Its output establishes only the listed checks,
not a proof of either original target.

\section{The claimed proof audit and the centered seven-term
theorem}\label{the-claimed-proof-audit-and-the-centered-seven-term-theorem}

This section records two different outcomes: an explicit counterexample
to a supporting proposition in a claimed proof, and a restricted
positive theorem in degree 20. Neither settles the Casas--Alvero
conjecture. The first is a proof audit; the second is an unpublished
mathematical result with exact certificates and internal adversarial
reviews. Priority, external refereeing, and publication are separate
matters.

\subsection{What the Ghosh audit
establishes}\label{what-the-ghosh-audit-establishes}

The audited source is Soham Ghosh, \emph{Proof of the Casas--Alvero
conjecture},
\href{https://arxiv.org/html/2501.09272v2}{arXiv:2501.09272v2}. The
arXiv history dates v2 to 21 March 2026. The captured manuscript has a
different internal date, 24 August 2026; these dates are recorded
separately. Its PDF SHA-256 is
\path{562050e1bff0fbfae2493951b4456485c61181381ba374271889d81d96bb5daa}.
The source identity, exact proof, replay, and dependency audit are in
\href{evidence/audit/README.md}{the audit evidence}.

For an ideal \(I\) in a ring \(R\), \(\mu_R(I)\) denotes its minimum
number of generators as an \(R\)-module. Over a local ring
\((R,\mathfrak m)\), this equals \(\dim_{R/\mathfrak m}I/\mathfrak mI\).
It is not the dimension or length of \(R/I\).

The paper defines \(I_n(j_1,\ldots,j_n)\) by applying root-recentering
automorphisms to elementary symmetric polynomials. Proposition 3.3
asserts equality between its global generator number and its generator
number at every minimal prime, provided the characteristic does not
divide \(\prod_{i=1}^n\binom n{i-1}\). Here \(n\) corresponds to CA
polynomial degree \(n+1\). For \(n=3\), that product is 9, so
characteristics 5 and 7 satisfy the stated hypothesis.

\subsubsection{An explicit counterexample in characteristic
5}\label{an-explicit-counterexample-in-characteristic-5}

For \((j_1,j_2,j_3)=(1,2,3)\), the definitions give, in \(R=k[x,y,z]\),

\[
F_1=-x(y-x)(z-x),\quad
F_2=xz-2xy-2yz+3y^2,\quad F_3=x+y-3z.
\]

Put \(I=(F_1,F_2,F_3)\). In characteristic 5, use invertible linear
coordinates

\[
\ell=x+y-3z,\qquad u=y-2z,\qquad t=z.
\]

Substitution gives \(F_2=(2t-2u)\ell+ut\) and

\[
F_1=-\ell^3-\ell^2u+\ell t^2+\ell tu-ut^2-u^2t+2u^3.
\]

Consequently \(I=(\ell,ut,u^3)\). At \(\mathfrak m=(\ell,u,t)\), these
monomial generators have independent classes in \(I/\mathfrak mI\). Thus
the local generator number is three; localization cannot increase the
number, and the presentation has three generators, so \(\mu_R(I)=3\).

The radical is the unique minimal prime \(\mathfrak p=(\ell,u)\). The
element \(t\) is a unit there, and

\[
I_{\mathfrak p}=(\ell,u)R_{\mathfrak p},\qquad
\mu_{R_{\mathfrak p}}(I_{\mathfrak p})=2.
\]

The last equality follows also from height two in this regular local
ring. This proves \(3\ne2\) under the proposition's hypotheses,
including over an algebraic closure.

The discrepancy is visible in

\[
I=(\ell,u)\cap(\ell,t,u^3).
\]

An embedded component carries nilpotent structure at the origin. More
directly,

\[
\mathfrak pI=(\ell^2,\ell u,u^2t,u^4).
\]

The class of \(u^3\) in \(I/\mathfrak pI\) is nonzero but killed by
\(t\). Localization removes the third generator; this torsion cannot be
cancelled.

The audit also reproduces the failure in the proof's literal translated,
augmented ideal. Set

\[
H_i=F_i(3s+A,s+B,3s+C),\quad L=A+B-3C,\quad U=B-2C.
\]

Then \(\mathfrak m^+=(L,U,B)\) and

\[
J^+=(H_1,H_2,H_3,B)=(L,B,Q,U^3),\qquad Q=3sU+2U^2.
\]

In \(k[s,L,U,B]/\mathfrak m^+J^+\), the class of \(U^3\) is nonzero:
setting \(L=B=0\) leaves the quotient by \((3sU^2+2U^3,U^4)\), whose
degree-three relation does not contain \(U^3\) alone. Yet

\[
3sU^3=U^2Q-2U^4\in\mathfrak m^+J^+.
\]

Thus the actual parameter is a zero divisor. The localization remains
nonzero after inverting \(s\); nonzero localization does not imply
absence of torsion.

\subsubsection{Characteristic 7 and the scope of the
refutation}\label{characteristic-7-and-the-scope-of-the-refutation}

In characteristic 7, \(\ell=x+y-3z,\ u=y-3z,\ t=z\) gives

\[
I=(\ell,u^2,ut^2)=(\ell,u)\cap(\ell,u^2,t^2).
\]

Again the global generator number is three and the number at the unique
minimal prime is two. Permuting \((x,y,z,0)\) and recentering gives all
24 distinct-entry ordered triples. They constitute one symmetry orbit,
not 24 unrelated mechanisms.

The same configuration does not refute the characteristic-zero
restriction. Eliminating \(F_3\) gives

\[
(y-3z)(y-2z)(2y-3z),\qquad5y^2-9yz+3z^2.
\]

Their affine resultant is \(315=3^2\cdot5\cdot7\), with no common
projective root at infinity in characteristic zero. The ideal then has
height three and both generator numbers are three. The established
conclusion is that Proposition 3.3 is false as stated and the proof
contains an invalid cancellation. This is neither a characteristic-zero
CA counterexample nor a disproof of Theorem A's conclusion.

\subsubsection{A separate characteristic-zero inference involving
nilpotents}\label{a-separate-characteristic-zero-inference-involving-nilpotents}

The \href{evidence/audit/DEPENDENCY_AUDIT.md}{dependency audit} traces
Proposition 3.3 through Corollary 3.9, Lemma 4.5, the filtered homology
argument, and the claimed descent in degree. It identifies another
invalid inference: a polynomial unit over an Artinian local ring can
have nonzero nilpotent coefficients. For example,

\[
A=\mathbb Q[\epsilon]/(\epsilon^2),\qquad
(1-\epsilon X)(1+\epsilon X)=1,
\]

although \(\epsilon\ne0\). This refutes the supporting coefficient
argument, not Lemma 4.5 itself for the special CA family.

A stronger abstract pressure test is retained. In
\(R=\mathbb Q[u,v,w]\), take the rows \((uv,u^2,v^2)\),
\((w^3,u^3,v^3)\), and their minors ideal \(D\). At
\(\mathfrak p=(u,v)\), \(D_{\mathfrak p}=(u^2,v^2)\). Thus \(uv\) is
nonzero and square-zero in the quotient, and \(w^3+uvX\) is a polynomial
unit. The corresponding polynomials \(uvX+w^3,\ u^2(X+u),\ v^2(X+v)\)
form a regular sequence: splitting the last two equations into the four
choices \(u=v=0\), \(u=0,X=-v\), \(v=0,X=-u\), or \(X=-u=-v\) shows that
their common zero set has dimension one in four variables. Surrounding
regular-sequence and Cohen--Macaulay properties therefore do not
automatically fix the inference. This model has a different degree
pattern and is explicitly not a CA configuration.

A valid cancellation repair needs the precise colon equality \((N:t)=N\)
for the relevant conormal module. A repair of the unit-polynomial step
needs control of nilpotents or a different argument. The audit does not
establish that such repairs are impossible. Later work on the general
repair target is reported separately.

\subsection{The positive sparse
theorem}\label{the-positive-sparse-theorem}

The authoritative statement is \href{evidence/seven-terms/PROOF.md}{the
seven-term proof}:

\textbf{Theorem.} Every nontrivial characteristic-zero degree-20 CA
polynomial has at least seven nonzero monomials after centering,
including the leading monomial.

Centering translates the unique root of \(f^{(19)}\) to zero. The CA
condition makes this a root of \(f\) too. After making the polynomial
monic, write

\[
f=X^{20}+\sum_{m\in S}c_mX^{20-m},\qquad
S\subseteq\{2,\ldots,19\},\quad c_m\ne0.
\]

The term count is \(1+|S|\). The index \(m\) is the exponent deficiency;
missing \(c_m\) means the Hasse derivative of order \(20-m\) vanishes at
zero. This is neither a distinct-root count nor a recycled-witness
count. It concerns the centered form, not every translate.

The last new exclusion is

\[
f=X^{20}+AX^{16}+BX^{15}+CX^{10}+DX^3+EX,\qquad ABCDE\ne0.
\tag{C}
\]

No such characteristic-zero polynomial is CA. Coefficient degenerations
are handled by preceding support results, not silently absorbed into
this exact-support statement.

\subsubsection{Complete support
coverage}\label{complete-support-coverage}

The inherited proof excludes at most five total terms. For exactly six,
the sieve examines all \(\binom{18}{5}=8568\) deficiency supports.
Established Lucas-visibility, one- and two-visible-coefficient
restrictions, placement restrictions, and the
Castryck--Laterveer--Ounaïes determinant criterion leave five:

{\def\LTcaptype{none} % do not increment counter
\begin{longtable}[]{@{}
  >{\raggedright\arraybackslash}p{(\linewidth - 4\tabcolsep) * \real{0.3333}}
  >{\raggedright\arraybackslash}p{(\linewidth - 4\tabcolsep) * \real{0.3333}}
  >{\raggedright\arraybackslash}p{(\linewidth - 4\tabcolsep) * \real{0.3333}}@{}}
\toprule\noalign{}
\begin{minipage}[b]{\linewidth}\raggedright
Label
\end{minipage} & \begin{minipage}[b]{\linewidth}\raggedright
Deficiency support
\end{minipage} & \begin{minipage}[b]{\linewidth}\raggedright
Exclusion
\end{minipage} \\
\midrule\noalign{}
\endhead
\bottomrule\noalign{}
\endlastfoot
A & \(\{3,4,10,18,19\}\) & Characteristic-13 seed, exact witness
collision, characteristic-zero resultants \\
B & \(\{3,10,16,17,19\}\) & Characteristic-13 identities including
coefficient degenerations \\
C & \(\{4,5,10,17,19\}\) & Six valuation branches below \\
D & \(\{4,10,12,17,19\}\) & Seed \(X^{20}+aX^{16}+cX^3+dX\) \\
E & \(\{8,10,16,17,19\}\) & Seed \(X^{20}+bX^4+cX^3+dX\) \\
\end{longtable}
}

The two-visible criterion is already in de Frutos Marín's thesis,
Proposition 3.5.5; the determinant criterion is CLO's. Neither is a new
method here. The complete frontier, its application of exclusions, and
the preceding lower bound are bundled in
\href{evidence/seven-terms/dependencies/casas-alvero-sixterm/PROOF.md}{the
predecessor proof}.

For A, the unique normalized nonmonomial seed is

\[
X^{20}+4X^{17}+X^{16}+4X^2+3X.
\]

Its Hasse orders 17, 16, and 2 have common witness 1, a simple
polynomial root, so their lifts coincide exactly. This gives

\[
X^{20}-1140X^{17}+14535X^{16}+tX^{10}
 +(-1589350-45t)X^2+(1575954+44t)X.
\]

Exact resultant certificates exclude this family, retaining \(t=0\)
separately. For B, \(X^{20}+aX^{17}+bX^4+cX^3+dX\) has no nonmonomial CA
point over an algebraic closure of characteristic 13. Its main-chart
certificate is

\[
U(u)C(u)R_3(u)+V(u)D(u)R_1(u)=(u+1)^{17}
\]

where \(u\ne-1\). Factors \(C,D\) retain zero-coefficient cases, and
separate certificates cover exceptional charts. This is not a search
restricted to prime-field rational points.

\subsubsection{Reduction of C and its six marked
cases}\label{reduction-of-c-and-its-six-marked-cases}

Use Hasse derivatives \(H_k(X^j)=\binom jkX^{j-k}\); in characteristic
zero they have the same roots as ordinary derivatives. Extend the
13-adic valuation and scale a root of least valuation to 1. All roots
and coefficients become integral, with a unit root retained. More
generally, in \(f=\sum\binom{20}{j}a_jX^{20-j}\), evaluation of each
monic normalized derivative at an integral common root inductively
proves integrality of every \(a_j\). This justifies the binomially
invisible coefficients disappearing in reduction.

For exact C, the order-10 witness \(z\) gives the stronger divisibility

\[
C=-184756z^{10}-8008Az^6-3003Bz^5\in13O.
\]

The coefficient \(A\) is a unit: its zero-residue seed charts either
contradict their derivative equations or force its exact witness to be
the simple mean root, implying \(A=0\). Thus its order-16 witness can be
scaled to 1, giving \(A=-4845\) and \(E=-1-A-B-C-D\).

The geometric seed classification gives exactly

\[
(\bar A,\bar B,\bar D,\bar E)
=(4,6,3,12),\ (4,6,2,0),\ (4,6,10,5).
\]

At the middle point a nonzero repeated root \(w\) must reduce to zero.
In

\[
f'(w)-f(w)/w
=19w^{19}+15Aw^{15}+14Bw^{14}+9Cw^9+2Dw^2,
\]

the last term has uniquely least valuation, a contradiction. This
includes its extension-field marked witnesses.

Let \(u,v,w\) witness orders 15, 3, and 1. Define

\[
h_3=X^{20}+4X^{16}+6X^{15}+3X^3+12X,\qquad
h_{10}=X^{20}+4X^{16}+6X^{15}+10X^3+5X.
\]

The remaining cases and exclusions are:

{\def\LTcaptype{none} % do not increment counter
\begin{longtable}[]{@{}
  >{\raggedright\arraybackslash}p{(\linewidth - 6\tabcolsep) * \real{0.2308}}
  >{\raggedright\arraybackslash}p{(\linewidth - 6\tabcolsep) * \real{0.2308}}
  >{\raggedleft\arraybackslash}p{(\linewidth - 6\tabcolsep) * \real{0.3077}}
  >{\raggedright\arraybackslash}p{(\linewidth - 6\tabcolsep) * \real{0.2308}}@{}}
\toprule\noalign{}
\begin{minipage}[b]{\linewidth}\raggedright
\((\bar v,\bar u,\bar w)\)
\end{minipage} & \begin{minipage}[b]{\linewidth}\raggedright
Forced equality
\end{minipage} & \begin{minipage}[b]{\linewidth}\raggedleft
\(\overline{C/13}\)
\end{minipage} & \begin{minipage}[b]{\linewidth}\raggedright
Final obstruction
\end{minipage} \\
\midrule\noalign{}
\endhead
\bottomrule\noalign{}
\endlastfoot
\((2,1,1)\) & \(u=w=1\) & 0 & Middle witness is exact 0, so \(C=0\) \\
\((2,1,4)\) & None initially & 12 & Coprime reduced polynomial and
divided middle derivative \\
\((2,2,1)\) & \(u=v,\ w=1\) & 3 & Second-precision contradiction at 4 \\
\((2,2,4)\) & \(u=v\) & 12 & Coprime reductions \\
\((11,1,3)\) & \(u=1\) & 1 & Coprime reductions \\
\((11,1,11)\) & \(u=1,\ v=w\) & 5 & Coprime reductions \\
\end{longtable}
}

Exact equalities follow from simple-root uniqueness or exhaustion of a
size-two cluster by an exact repeated root. Equality of residues alone
does not justify them.

\subsubsection{Why ramified lifts are
covered}\label{why-ramified-lifts-are-covered}

Two short principles underpin the proof. For integral equations in
positive-valuation deviations, an invertible residue Jacobian forces
each deviation to have at least the valuation of the constant defects.
Otherwise the invertible linear part has a minimum valuation strictly
below the constants and all nonlinear terms, so cancellation is
impossible. In the triple residue cluster at 4, the repeated-root
condition gives \(\nu(w-4)\ge1/2\): the constant and linear coefficients
of \(f'(4+(w-4))\) lie in \(13O\), while its quadratic coefficient is a
unit. Then \(f(w)=0\) forces \(\nu(f(4))>1\). This does not assume
\(w-4\in13O\).

If \(u=1\), then \(B=62016\). The Jacobians of \((f(v),H_3f(v))\) in
\((v,D)\) at \((2,3)\) and \((11,10)\) are

\[
\begin{pmatrix}9&6\\4&1\end{pmatrix},\qquad
\begin{pmatrix}0&7\\4&1\end{pmatrix},
\]

both with determinant 11. Hence write \(C=13k,\ v=v_0+13t,\ D=d_0+13l\).
The first equations reduce to

\[
9t+6l+8k+12=0,\quad4t+l+7k+6=0
\]

at \(v_0=2\), and

\[
7l+12k+2=0,\quad4t+l+6k+3=0
\]

at \(v_0=11\). The exact or clustered order-1 witness supplies a third
equation. The four solutions \((t,l,k)\) are
\((1,3,0),(5,7,12),(8,11,1),(0,6,5)\), respectively.

Dividing \(H_{10}f\) by 13 and a unit gives the reduction

\[
g_{9k}=X^{10}+11X^6+7X^5+9k.
\]

Exact Euclidean identities give

\[
\gcd(h_3,g_0)=X,\quad\gcd(h_3,g_4)=1,\quad
\gcd(h_{10},g_9)=\gcd(h_{10},g_6)=1.
\]

The first forces the middle witness to exact zero because zero is a
simple root of \(h_3\), whence \(C=0\). The others are immediate
contradictions.

For \((2,1,4)\), without an initial exact collision, put \(s=u-1\). The
order-15 equation gives

\[
B-62016=-15504s^2(10+10s+5s^2+s^3).
\]

For \(r=\nu(s)>0\), the unit Jacobian bounds \(v-2,D-3\) below by
\(\min(1,2r)\). In \((f(u)-f(1))/s\), the linear term in \(s\) has unit
coefficient 9 modulo 13. If \(r<1\), it has uniquely least valuation.
Hence \(r\ge1\), and the preceding \((2,4)\) computation applies.
Fractional ramification has been controlled rather than omitted.

For \(u=v\equiv2\), substitution gives the unit Jacobian
\(\begin{pmatrix}10&6\\4&1\end{pmatrix}\). With \(u=2+13r\), the first
equations are

\[
10r+6l+8k+7=0,\quad4r+l+7k+6=0.
\]

The \(w\equiv4\) row gives \(k=12\) and the same coprime pair. The
\(w=1\) row gives \((r,l,k)=(12,3,3)\), where \(\gcd(h_3,g_1)=X-4\) and
\(g_1'(4)=3\). Applying the unit-Jacobian bound to the exact equations
divided by 13 fixes the three parameters modulo \(13O\). The simple
divided-derivative root then forces \(z-4\in13O\). But direct expansion
gives \(\overline{f(4)/13}=6\ne0\), while \(f(z)-f(4)\in13^2O\). This
closes the last case.

\subsubsection{Arithmetic evidence and historical
corrections}\label{arithmetic-evidence-and-historical-corrections}

The combined replay checks 40 inherited items and ten new
normal/optimized runs, reconstructs six-row coverage, and includes an
altered-coefficient negative control. Finite-field certificates are
polynomial identities: interpolation uses more distinct extension-field
points than rigorously bounded degrees. Characteristic-zero resultant
coprimality has the required specialization and degree checks. The
compact C proof does not depend on the optional 1,544-determinant
resultant replay or exploratory collision logs. Internal reviews check
valuation implications separately from arithmetic.

The \href{evidence/AUDIT_NOVELTY_CORRECTION.md}{unchanged historical
correction} is essential. The early five-total-term bound follows
already from Massri's derivative-pair restrictions and CLO's mean and
determinant constraints. Its certificate remains correct; its assessment
as a potentially new theorem is withdrawn. Conversely, the correction's
then-current statement that no six-term theorem was claimed is
historical: the seven-term package supersedes it. The predecessor's
unresolved-C language is similarly superseded by the C proof, while its
certificates remain dependencies.

\subsection{Current status and priority
boundary}\label{current-status-and-priority-boundary}

The 23 September 2026 primary-source refresh still exposes Ghosh v2 as
the latest arXiv revision. This does not establish that no public
criticism or private repair exists. Ghosh is not the only author to have
claimed a full proof: \href{https://arxiv.org/abs/1707.04754v6}{Lu's
v6}, dated 16 March 2021, makes such a claim; it was not audited here.

Gasull's 2026 primer calls 24 the smallest open degree and cites
Draisma--de Jong. Its own proof covers degrees 4 and 5. The cited
theorem's degree-20 consequence was expressly withdrawn by its authors'
August 2011 erratum: the required characteristic-5 quartic exclusion
fails. Thus this citation chain supplies no new degree-20 proof.
Inheritance of an uncorrected status sentence is a plausible
explanation, not an established account of Gasull's reasoning.
\href{https://link.springer.com/article/10.1007/s44425-026-00047-6}{Primer},
\href{https://math-unibe.ch/jdraisma/publications/erratumcasasalvero.pdf}{erratum}.

The bounded sparse-priority search found no exact primary-source match
for the seven-total-term bound. This is not novelty clearance. The
ProofAtlas project page reports degree-20 certificate and Hensel work
without the full exponent systems and certificates needed to resolve
overlap. Its A1/A2/A3 labels must not be identified with this campaign's
A/B/C. A further thesis source remained unavailable. Standard valuation,
Hensel, determinant, and support methods are attributed to their
predecessors. The defensible result is an internally audited seven-term
theorem with unresolved historical priority, alongside a reproducible
proof-gap audit. Unrestricted degree 20 and the all-degree conjecture
remain unproved by this work.

\section{The unrestricted degree-20
investigation}\label{the-unrestricted-degree-20-investigation}

This section concerns a hypothetical nontrivial degree-20 Casas--Alvero
polynomial in characteristic zero. No sparsity assumption is imposed
unless a particular support is named. The results comprise complete
residue classifications, one whole residue-branch exclusion, several
lifting reductions, conditional exclusions, and finite algebraic targets
that remain unevaluated. They do not prove the degree-20 conjecture.
Eight of the nine characteristic-17 branches below remain unresolved as
whole branches.

\subsection{Normalization, integrality, and the simple
mean}\label{normalization-integrality-and-the-simple-mean}

Write \(H_kf=f^{(k)}/k!\) for the \(k\)-th Hasse derivative. Translate
the root shared with \(H_{19}f\) to zero and make \(f\) monic. Its mean
root is now zero. At a prime \(p\), scale by a nonzero root of minimum
valuation. All roots are integral and at least one is a unit. Write

\[
 f(X)=\sum_{j=0}^{20}C_ja_jX^{20-j},\quad C_j=\binom{20}{j},
 \qquad a_0=1,\quad a_1=a_{20}=0,
\] \[
 G_j(X)=\frac{H_{20-j}f(X)}{C_j}
       =\sum_{i=0}^{j}\binom ji a_iX^{j-i}.                 \tag{D1}
\]

Every \(G_j\) shares a root \(w_j\) with \(f\). Induction in \(j\),
using the coefficient one on \(a_j\) in \(G_j(w_j)=0\), proves that
\textbf{all normalized coefficients are integral}. Ordinary coefficient
integrality alone would not prove this. The residue arguments retain
these normalized coefficients even when \(C_j\) vanishes modulo \(p\).

Algebraic counterexamples suffice for an existence argument. The root
and derivative incidences, with a normalization enforcing two distinct
roots, are polynomial equations over the rationals. A complex point
would give an algebraic point by the weak Nullstellensatz. Thus finite
extensions of local fields suffice. No argument below assumes these
extensions unramified unless that conclusion is proved.

The exact mean is simple. If it were multiple, \(a_{19}=0\). Normalize
integrally at 19. Every \(C_j\), \(2\le j\le18\), is divisible by 19;
the other nonleading coefficients vanish by centering, the root at zero,
and \(a_{19}=0\). The reduction would be \(X^{20}\), contradicting the
retained unit root. This proves the needed special case of the cited
simple-mean theorem directly.

A separated-cluster lemma will be used repeatedly. Suppose a cluster
contains \(m\) roots, with multiplicity, and common witnesses for Hasse
orders \(1,\ldots,m-1\). If degree-\(m\) Hasse--CA holds in the residue
characteristic, the cluster is one exact root. Otherwise translate a
cluster root to zero and scale by a nonzero difference of least
valuation. The cluster factor reduces to a degree-\(m\) polynomial with
at least two roots; the outside factor reduces to a nonzero constant.
Hasse's product rule transfers all the indicated incidences to the
cluster reduction, a contradiction. Actual witnesses and the full
separated cluster are essential; occupancy counts alone do not suffice.

In characteristic \(p\), this applies when \(m=p^e\): a highest
nonconstant term below \(X^{p^e}\), if present, would have a nonzero
constant Hasse derivative of its own order. It also applies to degrees
2,3,4 in characteristic 17. The cubic reduces after centering to
\(X^3+bX\); its first-derivative common-root equations force \(b=0\).
The centered quartic is \(X^4+aX^2+bX\). If \(a=0\), then \(b=0\);
otherwise normalize an \(H_2\) witness to one, obtaining
\(X^4-6X^2+5X\), coprime to \(4X^3-12X+5\) over \(\mathbb F_{17}\).

Proof sources:
\href{evidence/full/LIFT_CONSEQUENCES_17.md}{normalization and simple
mean}, \href{evidence/full/two_adic/CLUSTER_COLLAPSE.md}{cluster
collapse}, and
\href{evidence/full/wild17/blowup/m4/check_lift_consequences.py}{the
quartic arithmetic check}.

\subsection{The complete reductions at 2 and
5}\label{the-complete-reductions-at-2-and-5}

At 2, Lucas's rule gives \[
 \bar f=X^{20}+\bar a_4X^{16}+\bar a_{16}X^4.
\] Scale a unit root to one, so \(\bar a_4+\bar a_{16}=1\). An \(H_4\)
common root satisfies \(r^{16}=\bar a_{16}\); substitution in \(\bar f\)
gives \(\bar a_4\bar a_{16}=0\). There are exactly two ordinary seeds:

{\def\LTcaptype{none} % do not increment counter
\begin{longtable}[]{@{}lll@{}}
\toprule\noalign{}
Type & Seed & Cluster sizes at \(0,1\) \\
\midrule\noalign{}
\endhead
\bottomrule\noalign{}
\endlastfoot
A & \(X^{20}+X^{16}\) & \(16,4\) \\
B & \(X^{20}+X^4\) & \(4,16\) \\
\end{longtable}
}

Both satisfy every residue Hasse--CA condition. Every witness reduces to
zero or one. Its equation determines \(\bar a_j=0\) when \(\bar w_j=0\),
and \(\bar a_j=\sum_{i<j}\binom ji\bar a_i\) otherwise. This proves
containment in \(\mathbb F_2\) over any residue-field extension. Degrees
4 and 16 are forced, leaving 16 witness bits per type. Complete
enumeration gives 65,536 marked assignments in each type, representing
873 and 1,128 distinct coefficient patterns.

A genuine integral lifting cut improves this. Put \((m,k)=(4,16)\) in A
and \((16,4)\) in B, and scale the actual unit \(G_m\) witness exactly
to one. Set \(F=f(1)\), \(P=f(w_k)\), \(g=G_m(1)\), \(q=G_k(w_k)\). The
integer polynomial \[
 S=P+w_k^mq+(g+1)(F+g+q)                                  \tag{D2}
\] has every coefficient even: modulo 2 substitute \(F=1+a_m+a_k\),
\(P=w_k^{20}+a_mw_k^k+a_kw_k^m\), \(g=1+a_m\), \(q=w_k^k+a_k\), and use
\(m+k=20\). Thus \(S/2\) is integral and vanishes at every
characteristic-zero incidence point. Its reduction gives \[
 \bar a_8+\bar a_{12}+\bar a_{18}=0\quad(A),\qquad
 \bar a_2+\bar a_{12}+\bar a_{18}=0\quad(B).                \tag{D3}
\] No merely positive-valuation coordinate difference is divided by 2.
The identity therefore permits arbitrary ramification.

The cuts leave 465 and 603 coefficient patterns, respectively, and
40,960 marked assignments per type. Cluster collapse and the simple mean
remove another 4 marked assignments in A and 4,096 in B: the zero
cluster cannot contain all witnesses through Hasse order 15 in A, or
through order 3 in B. The resulting total is \textbf{77,820 marked
assignments and 1,068 coefficient patterns}. No additional entire
coefficient pattern is removed by this occupancy test.

The active-witness equations also give \(a_m+1\in2O\), \(a_k\in2O\), and
\(f\equiv X^{20}-X^k\pmod{2O[X]}\). Comparing coefficient valuations at
a zero-cluster root gives \[
 \nu_2(r)\ge1/14\quad(A),\qquad \nu_2(r)\ge1/2\quad(B)
\] for nonzero \(r\), and then \(\nu_2(a_{16})\ge8/7\) in A and
\(\nu_2(a_4)\ge2\) in B. These are exact rational comparisons of
binomial valuations with exponent gaps. Neither type is excluded.
Low-precision incidence points exist modulo 16 in A and modulo 4 in B;
they are not characteristic-zero lifts. In A one such point is
\(a_4=-1,a_{19}=3875\), all other coefficients zero, with witnesses
\(w_4=w_{19}=1\) and all others zero. In B use \(a_{16}=-1\), all others
zero, with \(w_{16}=1\). The displayed ring precisions satisfy all
incidences but do not assert infinite liftability. The
\href{evidence/full/two_adic/REDUCTION_AND_LIFTING.md}{two-adic proof
and replays} include both independent censuses and further conditional
precision facts.

At 5, the reduction is \[
 \bar f=X^{20}+aX^{15}+bX^{10}+cX^5=q(X)^5,\qquad
 q=X^4+AX^3+BX^2+CX.
\] Frobenius and the Hasse identities make \(q\) a characteristic-5
Hasse--CA quartic. Translate its own \(H_3\) root to zero for this
classification, giving \(Q=X^4+UX^2+VX\). If \(U=0\), the \(H_1\)
condition forces \(V=0\), contrary to two distinct roots. Otherwise
normalize an \(H_2\) witness to one. Since \(H_2Q=X^2+U\), its equation
and \(Q(1)=0\) give \(U=-1,V=0\). The original root zero is therefore
either the double root or a simple root. Up to scaling preserving zero,
the two seeds are \[
 X^{20}-X^{10},\qquad X^{20}+X^{15}+3X^5,
\] with cluster sizes \(10,5,5\) at \(0,1,-1\), or \(5,10,5\) at
\(0,1,2\). Both satisfy all residue incidences. The internal translation
of \(q\) is not an arbitrary translation of the centered
characteristic-zero \(f\). This is a full-support reduction, not a
contradiction; see the
\href{evidence/full/five_adic/REDUCTION.md}{five-adic proof}.

\subsection{Nine characteristic-17 seeds and one whole-branch
exclusion}\label{nine-characteristic-17-seeds-and-one-whole-branch-exclusion}

At 17 the visible polynomial is \[
 h=X^{20}+aX^{18}+bX^{17}+cX^3+dX^2+eX.                  \tag{D4}
\] Here \(a,b,c,d,e\) are ordinary residue coefficients. Over every
algebraically closed characteristic-17 field, the nonmonomial Hasse--CA
polynomials of this form have the following representatives after the
indicated nonzero scaling.

{\def\LTcaptype{none} % do not increment counter
\begin{longtable}[]{@{}
  >{\raggedright\arraybackslash}p{(\linewidth - 12\tabcolsep) * \real{0.1154}}
  >{\raggedleft\arraybackslash}p{(\linewidth - 12\tabcolsep) * \real{0.1538}}
  >{\raggedleft\arraybackslash}p{(\linewidth - 12\tabcolsep) * \real{0.1538}}
  >{\raggedleft\arraybackslash}p{(\linewidth - 12\tabcolsep) * \real{0.1538}}
  >{\raggedleft\arraybackslash}p{(\linewidth - 12\tabcolsep) * \real{0.1538}}
  >{\raggedleft\arraybackslash}p{(\linewidth - 12\tabcolsep) * \real{0.1538}}
  >{\raggedright\arraybackslash}p{(\linewidth - 12\tabcolsep) * \real{0.1154}}@{}}
\toprule\noalign{}
\begin{minipage}[b]{\linewidth}\raggedright
Row
\end{minipage} & \begin{minipage}[b]{\linewidth}\raggedleft
\(a\)
\end{minipage} & \begin{minipage}[b]{\linewidth}\raggedleft
\(b\)
\end{minipage} & \begin{minipage}[b]{\linewidth}\raggedleft
\(c\)
\end{minipage} & \begin{minipage}[b]{\linewidth}\raggedleft
\(d\)
\end{minipage} & \begin{minipage}[b]{\linewidth}\raggedleft
\(e\)
\end{minipage} & \begin{minipage}[b]{\linewidth}\raggedright
Witness normalized to 1
\end{minipage} \\
\midrule\noalign{}
\endhead
\bottomrule\noalign{}
\endlastfoot
1 & 0 & 0 & 16 & 0 & 0 & \(H_3\) \\
2 & 14 & 0 & 16 & 0 & 3 & \(H_3\) \\
3 & 14 & 8 & 16 & 12 & 0 & \(H_3\) \\
4 & 14 & 0 & 0 & 11 & 8 & \(H_{18}\) \\
5 & 14 & 2 & 0 & 0 & 0 & \(H_{18}\) \\
6 & 14 & 2 & 0 & 11 & 6 & \(H_{18}\) \\
7 & 14 & 2 & 0 & 14 & 3 & \(H_{18}\) \\
8 & 0 & 16 & 0 & 0 & 0 & \(H_{17}\) \\
9 & 0 & 16 & 0 & 14 & 3 & \(H_{17}\) \\
\end{longtable}
}

Containment in \(\mathbb F_{17}\) is a conclusion, not an initial
restriction. Completeness uses explicit elimination certificates and
elementary boundaries. Put \[
 U=X^{17}+c,\quad V=X^3+aX+b,\quad
 Q=dX^2+(e-ac)X-bc.
\] Then \(h=UV+Q\), \(H_3h=U\), \(H_{17}h=V\), \(H_2h=3XU+d\),
\(H_{18}h=3X^2+a\), and \(H_1h=(3X^2+a)U+2dX+e-ac\). Other derivative
conditions hold at zero. If \(c\ne0\), normalize the unique \(H_3\) root
to one, so \(c=-1,e=-a-b-d\) and \(Q=(X-1)(dX-b)\). An \(H_{17}\) common
root then gives three covering charts: \(b=0\); witness one; or witness
\(u\ne1\), with \(a=-u^2-d,b=ud\). The remaining conditions are the
three resultants with \(H_{18},H_2,H_1\).

In the first chart their ideal contains \(d^{20}\) and a polynomial
specializing at \(d=0\) to \(a^3(a+3)^{17}\). The second chart has a
certificate for 1. In the third, with \(s=u,t=d\), certificates give
\(t^{20}(t+5)\), \(t^{20}(s+5)\), and a polynomial specializing at
\(t=0\) to \(s^6(s^2-3)^{17}\). These force rows 1--3. If \(c=0,a\ne0\),
normalize \(a=-3\), with parameters \(b=s,d=t\). The resultant ideal
contains \[
 [s(s-2)]^{19},\quad[st+6s-2t+5]^{19},\quad
 [t(t-11)(t-14)]^{18},
\] giving rows 4--7. If \(c=a=0,b\ne0\), normalize \(b=-1,e=-d\). For
\(d\ne0\), an \(H_2\) witness satisfies \(-r^{16}(r-1)^2(2r+1)=0\); the
possible coefficient values give row 9 or fail \(H_1\). The case \(d=0\)
is row 8. If \(a=b=c=0,d\ne0\), normalization gives \(d=-3,e=2\); the
remaining equations force \(r=7\), but \(h(7)/7=1\). The final binomial
case fails unless it is monomial.

The
\href{evidence/full/support_frontier/prime17/CLASSIFICATION.md}{classification
appendix} contains the membership certificates, checked by coefficient
multiplication. Twelve resultants are independently reconstructed using
12,895 exact evaluations and proved interpolation bounds over a
sufficiently large extension field. Every retained seed is checked by
univariate gcds. This is an algebraic-closure classification, not a
prime-field search.

\textbf{Row 3 cannot lift.} Its \(\gcd(h,H_1h)=X\), while zero has
multiplicity exactly two. An exact repeated root \(r\) of a lift must
reduce to zero, but the simple-mean theorem says \(r\ne0\). Then
\(X(X-r)^2\mid f\), and reduction of its integral monic factorization
forces \(X^3\mid h\), a contradiction. This excludes the entire row,
including invisible coefficients and ramified lifts.

The other exact consequences used below are \[
\begin{array}{c|l}
4&a_2=-1,\ a_3=a_{17}=0,\\
6,7&a_2=-1,\ a_3=2,\ a_{17}=0,\\
9&a_2=a_{17}=0,\ a_3=-1,\ (X-1)^3\mid f.
\end{array}                                                    \tag{D5}
\] For rows 4,6,7, simple residue classes force equality of the relevant
exact witnesses; \(G_2(1)\) and \(G_3(1)\) give the constants. In row 9
the gcds for \(H_1,H_2,H_{17}\) are \((X-1)^2,X-1,X-1\). The three-root
cluster has its \(H_1,H_2\) witnesses and collapses exactly by the
degree-three lemma. The simple mean gives \(a_2=a_{17}=0\);
normalization at its \(H_{17}\) witness gives \(a_3=-1\). See the
\href{evidence/full/LIFT_CONSEQUENCES_17.md}{lifting consequences}.

\subsection{Complete local descriptions for rows
4,6,7,9}\label{complete-local-descriptions-for-rows-4679}

These reductions impose a square system with nonsingular residue
Jacobian, or a controlled quadratic exception, leaving one exact scalar
equation. Solving the square system does not assert that last equation.
Uniqueness also covers ramification: if two integral solutions with the
same residues differ at minimum valuation \(\gamma>0\), an invertible
integral Jacobian preserves that minimum in the linear Taylor term;
nonlinear terms have value at least \(2\gamma\), so cannot cancel it.

Write \(u_j=a_j\), \(4\le j\le16\). In row 9, solving only
\(f(1)=f'(1)=0\) gives \[
\begin{split}
 f_u={}&X^{20}-1140X^{17}+\sum C_ju_jX^{20-j}\\
 &+\left(18221-\sum(19-j)C_ju_j\right)X^2\\
 &+\left(-17082+\sum(18-j)C_ju_j\right)X.
\end{split}                                                    \tag{D6}
\] Sums in this subsection run over the middle degrees. The residue is
the fixed seed for every integral \(u\): each parameter derivative is
\(C_j[X^{20-j}-(19-j)X^2+(18-j)X]\), divisible by 17. Mark each middle
witness residue. Use exact zero or one for those classes, justified by
the simple mean and triple collapse; for other classes introduce root
variables with equations \(f_u(r_j)=0\). Add all \(G_j(w_j)=0\). The
reduced Jacobian has blocks \[
 \begin{pmatrix}D&0\\ B&L\end{pmatrix},
\] where \(D\) has nonzero simple-root derivatives on its diagonal and
\(L\) is lower triangular with diagonal one. The triangular equations
put all coefficient residues in the seed's finite splitting field. Every
marking thus has a unique unramified square-system lift, also unique in
ramified ambient fields. The missing triple equation is \[
 T(u)=-8037+\sum_{j=4}^{16}\binom{19-j}{2}\frac{C_j}{17}u_j=0,
 \qquad H_2f_u(1)=17T(u).                                    \tag{D7}
\] There are 18 distinct residue roots, with splitting field
\(\mathbb F_{17^{10}}\), giving naive marking bound \(18^{13}\).
Prime-field witnesses alone would omit cases.

Row 6 has its \(H_2\) witness \(s\) reducing to 6 and its exact double
root \(r\) reducing to 10. Eliminate \[
 A=a_{18}=-s^{18}+153s^{16}-1632s^{15}
       -\sum\binom{18}{j}u_js^{18-j},\qquad
 20a_{19}=-2091-\sum C_ju_j-190A.
\] Use \(f_{u,s}(s)=0\), the critical equation \(G_{19}(r)=0\), the
other simple-root equations, and all middle equations. The root block is
lower triangular with nonzero diagonal \(7,9\), followed by simple-root
derivatives. The middle block has diagonal one. Here 7 is the total
\(s\)-derivative, including the eliminated coefficients. The remaining
equation is \(f_{u,s}(r)=0\).

In row 7 the \(H_2\) witness is the simple exact root one. Solving
\(f(1)=H_2f(1)=0\) gives \[
 190a_{18}=-281200-\sum\binom{20-j}{2}C_ju_j,
\] \[
 20a_{19}=279109+\sum\left(\binom{20-j}{2}-1\right)C_ju_j.
\] The double cluster reduces to 7. Its critical equation
\(G_{19}(r)=0\) has derivative 11 modulo 17, and parameter derivatives
vanish modulo 17. The same block argument applies, leaving \(f_u(r)=0\)
untested. Both reductions cover every marking and arbitrary ramified
candidates, but exclude neither row.

Row 4 has two double residue clusters. Its seed factors as
\(X(X+7)(X+11)(X+16)(X^2+3X+3)^2\) times three irreducible quartics, so
it splits over \(\mathbb F_{17^4}\). Let \(\alpha,\beta\) be the
quadratic roots and choose which supplies the repeated root. The \(H_2\)
witness \(s\) reduces to 6. Eliminate \[
 a_{18}=-s^{18}+153s^{16}-\sum\binom{18}{j}u_js^{18-j},\qquad
 20a_{19}=189-\sum C_ju_j-190a_{18}.
\] The total derivative of \(f_{u,s}(s)\) is 9 modulo 17, giving an
integral analytic solution \(s=S(u)\). Simple classes give analytic root
functions; \(G_{19}(R_\alpha)=0\) gives a critical-root function with
derivative 2. The separated beta factor is \((X-c(u))^2-d(u)\), with
\(d(u)\) coefficientwise divisible by 17. For each witness in this
cluster choose independently \(c(u)+z\) or \(c(u)-z\). The middle
Jacobian in \(u\) has unit triangular diagonal, so \(u=u(z)\) is an
integral power-series vector. The remaining cluster equation
\(z^2-d(u(z))=0\) is, by Weierstrass preparation, a unit times a monic
quadratic congruent to \(z^2\) modulo 17.

Every actual row-4 point, with all its roots, therefore lies in an
extension of degree at most two over the unramified degree-four
extension of \(\mathbb Q_{17}\). Its local degree is at most eight and
its ramification index at most two. Signed charts retain both roots of
the second cluster. The additional scalar equation is still
\(f_{u(z)}(R_\alpha(u(z)))=0\); all charts have not been excluded.

One complete row-4 support, with middle set \(\{4,10,12\}\), is
excluded. Its divided critical value leaves only two conjugate
orientations among \(2\cdot17^3=9,826\) complete markings, with
witnesses \((10,1,6)\) and coefficients \((1,4,9)\). Their square-system
Jacobian is
\(\left(\begin{smallmatrix}2&12\\0&9\end{smallmatrix}\right)\). The
unique lift has critical value \(17^2(9+4\alpha)\pmod{17^3}\), nonzero
because \(\alpha^2+3\alpha+3=0\) is irreducible. This excludes that
support in ramified extensions too, not all row 4.

Detailed equations, factorization checks and separate reviews are in
\href{evidence/full/tame17/ROW4_QUADRATIC_LIFTING.md}{the row-4 lifting
proof},
\href{evidence/full/tame17/ROW4_SMALLEST_SUPPORT_EXCLUSION.md}{its
support exclusion}, and the
\href{evidence/full/tame17/ROW6_ETALE_AUDIT.md}{row-6},
\href{evidence/full/tame17/ROW7_ETALE_AUDIT.md}{row-7}, and
\href{evidence/full/tame17/ROW9_ETALE_AUDIT.md}{row-9} lifting audits.

\subsection{Complete row-9 support calculations and the unevaluated
norm}\label{complete-row-9-support-calculations-and-the-unevaluated-norm}

Canonical zero-witness choices and the audited support sieve reduce row
9 to 240 eligible coefficient-support systems. An exactly active
coefficient cannot have a witness reducing to zero: that simple class
contains only the exact mean. Its coefficient residue may nevertheless
be zero; these degenerations remain included. Each active middle witness
ranges over the 17 nonzero residue roots, including extension-field
values. The triangular coefficient recursion is followed by the residue
test for (D7). Every survivor has a unique square-system lift; a nonzero
exact digit of \(T\) excludes it. Frobenius-orbit grouping does not omit
conjugate or extension-field markings.

The computation uses a slightly different square subsystem from the
genuine-candidate description above. Put \(q_u=f_u/[X(X-1)^2]\). After
the triangular derivative recursion, impose \(q_u(x_r)=0\) for each
distinct selected residue root. Its reduced Jacobian is diagonal with
entries \(\overline q'(r)\ne0\), because all parameter dependence of
\(q_u\) is divisible by 17. In this subsystem the root near residue one
is allowed to move; it is not fixed equal to one before \(T=0\) is
imposed. Every genuine candidate is covered. Unit-Jacobian uniqueness
over ramified ambient extensions shows that an exact solution agrees
with a saved approximation modulo \(17^N\). Thus a saved value of \(T\)
of valuation less than \(N\) excludes that marking. This is the lifting
implication used by all three batches.

The support-sieve dependency is explicit. The frozen full inventory has
2,482 remaining centered supports across term counts 7 through 17, after
Lucas visibility, the stated prior-art conditions, and the separately
proved sparse exclusions. It is a list of necessary supports, not CA
polynomials. The row-9 exact conditions select the 240 systems from that
inventory. Subsequent local exclusions are recorded separately rather
than silently changing the frozen count; see
\href{evidence/full/support_frontier/REPORT.md}{the support diagnostic}.

The three completed batches exclude all \textbf{79 systems with at most
six active middle coefficients}, among the 240 systems in the cover. The
remaining \textbf{161 systems have not been executed}. The complete
searches cover \textbf{1,261,339,055 marked residue assignments}. Their
819 first-residue survivors form 307 Frobenius orbits, all excluded at
precision at most \(17^4\): 811 survivor markings are excluded at
\(17^2\), six at \(17^3\), and two at \(17^4\). Assignments failing the
first-residue test need no lift.

{\def\LTcaptype{none} % do not increment counter
\begin{longtable}[]{@{}
  >{\raggedright\arraybackslash}p{(\linewidth - 8\tabcolsep) * \real{0.1579}}
  >{\raggedleft\arraybackslash}p{(\linewidth - 8\tabcolsep) * \real{0.2105}}
  >{\raggedleft\arraybackslash}p{(\linewidth - 8\tabcolsep) * \real{0.2105}}
  >{\raggedleft\arraybackslash}p{(\linewidth - 8\tabcolsep) * \real{0.2105}}
  >{\raggedleft\arraybackslash}p{(\linewidth - 8\tabcolsep) * \real{0.2105}}@{}}
\toprule\noalign{}
\begin{minipage}[b]{\linewidth}\raggedright
Active middle coefficients
\end{minipage} & \begin{minipage}[b]{\linewidth}\raggedleft
Systems
\end{minipage} & \begin{minipage}[b]{\linewidth}\raggedleft
Marked assignments
\end{minipage} & \begin{minipage}[b]{\linewidth}\raggedleft
Residue survivors
\end{minipage} & \begin{minipage}[b]{\linewidth}\raggedleft
Frobenius orbits
\end{minipage} \\
\midrule\noalign{}
\endhead
\bottomrule\noalign{}
\endlastfoot
Three or four & 7 & 506,039 & 20 & 7 \\
Five & 21 & 29,816,997 & 113 & 48 \\
Six & 51 & 1,231,016,019 & 686 & 252 \\
Total & 79 & 1,261,339,055 & 819 & 307 \\
\end{longtable}
}

The first two batches have complete FLINT enumerations and separate
native enumerations, followed by independent survivor and exact-lift
checks. For the final batch, the guarded native producer completed 46
systems before its wall limit; a continuation ran only the remaining
five. A fresh complete native run agreed byte for byte. This is a repeat
of the same enumeration algorithm, \textbf{not an independent exhaustive
FLINT replay}. Separate code checks verify the support inventory,
complete enumeration blocks, surviving markings, Frobenius coverage, and
lifting arithmetic; normal and optimized runs agree. The final batch was
the last search attempted for this review.

The complete proofs and receipts are in
\href{evidence/full/collective17/exclusions/}{the batch evidence}. The
separate prime-field-witness subcase follows from the cited Massri
theorem after an affine coordinate change; see
\href{evidence/full/collective17/prime-field/PRIOR_ART.md}{its prior-art
argument}. That external computational proof was not replayed and is not
presented as a new exclusion. Its application has a short exact bridge:
row 9 has only the prime-field residue roots \(0,1,-2\). Their
genuine-candidate clusters are exact zero, exact triple one, and a
unique simple root \(r\). Active witnesses must therefore be \(1\) or
\(r\); inactive derivatives can use zero. All incidences use three exact
roots. The change \(F(x)=f(1-x)\) gives precisely Massri's degree-20
three-recycled-root configuration, with repeated root zero and mean one.

There is also a single exact collective algebraic target. Set
\(g_u=f_u/(X-1)^2\), a monic degree-18 integer polynomial, so
\(g_u(1)=17T(u)\). Define \[
 R_j(u)=\operatorname{Res}_X(g_u,G_j),\qquad
 B=\mathbb Z_{17}\langle u_4,\ldots,u_{16}\rangle/(R_4,\ldots,R_{16}).
                                                               \tag{D8}
\] The brackets denote the 17-adically completed polynomial ring. When
\(T=0\), \(g_u\) and \(f_u\) have the same roots. Thus
\(R_4=\cdots=R_{16}=T=0\) is equivalent to the complete incidence system
in this normalized integral branch. Without \(T=0\), using \(g_u\) is an
enlargement, not that equivalence.

Modulo 17, \(g_u\) is the fixed squarefree degree-18 polynomial
\(D=h/(X-1)^2\). Each \(\bar R_j=\operatorname{Res}(D,\bar G_j)\) is
monic of degree 18 in the fresh variable \(u_j\) and has no later
\(u_i\). Successive monic division gives residue basis
\(\prod u_j^{e_j}\), \(0\le e_j<18\); these reductions form a regular
sequence. The flatness lift is direct. If a complete 17-torsion-free
ring \(A\) has an element \(r\) with nonzerodivisor reduction and
\(17x=ry\), then \(y=17z\) after reduction; cancellation gives \(x=rz\).
Thus \(A/(r)\) stays 17-torsion-free. Apply this successively to the
resultants. Completeness lifts the residue basis to generators; repeated
division of any relation by 17 proves independence. Consequently \(B\)
is finite free of rank \[
 18^{13}=20,822,964,865,671,168.
\]

Let \(N\) be the determinant of multiplication by \(T\) on \(B\). The
whole integral row-9 branch is empty in characteristic zero \textbf{if
and only if \(N\ne0\) in \(\mathbb Q_{17}\)}. Indeed this is equivalent
to \(T\) being a unit in the finite generic algebra and therefore to its
quotient by \(T\) being zero. A nonzero quotient has a geometric point;
its coordinates are integral because \(B\) is finite over
\(\mathbb Z_{17}\). Nilpotents and coincident markings are retained.
Nonvanishing is sufficient; \(N\) need not be a 17-adic unit and could
be highly divisible by 17.

For a canonical middle support \(J\), use \(g_u/X\) and only \(j\in J\).
The exact identity \(R_j=u_j\operatorname{Res}(g_u/X,G_j)\) justifies
deleting the mean for active witnesses. The resulting completed algebra
has rank \(17^{|J|}\). Across the 240 supports these ranks sum to
813,975,725,115,600. Neither these generic norms nor \(N\) have been
evaluated. The exact resultants need not be triangular over the
integers: the proof used their triangular reductions and completion.

A compact circuit uses degree-17 or degree-18 companion matrices to
represent each resultant. It does not construct the enormous final
multiplication matrix. There is also a small Frobenius description of
the residue domain: \[
 (X-1)D=X^{17}(X^2+X+1)-3X.
\] Every root \(r\) obeys \(r^{17}=3r/(r^2+r+1)\), with nonzero
denominator. For \(r\ne1\), the change \(y=7(r+1)/(r-1)\) transforms
this to \(y^{17}=y^2-5\); the root \(r=1\) corresponds to infinity. This
description has not supplied a collective norm evaluation. The
\href{evidence/full/collective17/elimination/FINITE_FLAT_REDUCTION.md}{finite-free
proof and presentation} give the coefficient identities, support-list
dependency and replay.

\subsection{Wild row 8: scales and quartic first
models}\label{wild-row-8-scales-and-quartic-first-models}

Here \(\bar f=X^{17}(X^3-1)\). Normalize its actual \(G_3\) unit witness
to one, giving \(a_3=-1-3a_2\). Let \(\delta>0\) be the least valuation
of a nonzero root in the 17-root zero cluster. The witnesses of
\(G_2,G_{17},G_{18},G_{19}\) are in that cluster. Their equations give
\[
 \nu(a_2)\ge2\delta,\qquad
 \nu(a_j)\ge\min(j\delta,1+(j-16)\delta),\quad j=17,18,19.
\] At a root attaining \(\delta\), the \(X^{17}\) term has value
\(17\delta\); all others have value at least
\(\min(20\delta,1+4\delta)\). Cancellation requires \(\delta\ge1/13\).
In particular the final three coefficient valuations are at least
\(1+\delta,1+2\delta,1+3\delta\). The exact equation \[
 0=-17\cdot67-17\cdot190a_2+\sum_{j=4}^{16}C_ja_j
        +1140a_{17}+190a_{18}+20a_{19}
\] may consequently be divided by 17 and reduced, yielding \[
 \sum_{j=4}^{16}(C_j/17)\bar a_j=-1.                       \tag{D9}
\] The stronger valuation estimates, rather than residue divisibility
alone, justify this division with arbitrary ramification.

All middle witness residues lie in \(\{0,1,\zeta,\zeta^2\}\), where
\(\zeta\) is a primitive cube root in \(\mathbb F_{17^2}\). Triangular
equations determine every normalized coefficient residue. Exactly
233,310 of the \(4^{13}=67,108,864\) marked assignments satisfy (D9).
Two exact common-root obstructions leave 180,341: no root is common to
\(f,H_4f,H_{16}f\), or to \(f,H_5f,H_{10}f,H_{15}f\). Translate such a
root to zero and normalize integrally at 2 or 5. The specified exact
coefficient zeros would remove every nonleading Lucas-visible term,
contradicting a unit root. Simplicity of the three unit residue classes
permits these obstructions on equal unit labels. It does not identify
arbitrary zero-cluster witnesses. A single unit middle witness, if
present, must have normalized degree 11, residue one, and coefficient
residue 11; that case is not excluded.

For a remaining marking let \(J\) be the largest middle degree with a
unit witness and \(L\) the largest with nonzero coefficient residue.
Then \(4\le L\le J\le16\). Zero-witness recursion for \(j>J\) gives
\(\nu(a_j)\ge(j-J)\delta\), and the final three recurrences give
\(\nu(a_j)\ge\min(j\delta,1+(j-J)\delta)\). The minimum-term argument
improves to \(\delta\ge1/(J-3)\). Conversely, factor \(f=UV\), with
\(U\) its degree-17 zero-cluster factor and \(V\) its integral
degree-three unit factor. The coefficient of \(X^{20-L}\) has valuation
at least \((L-3)\delta\), but is \(C_La_L\) of valuation one. Hence \[
 \frac1{J-3}\le\delta\le\frac1{L-3}.                      \tag{D10}
\] Exactly 179,765 markings have \(J=L\), fixing the scale. The other
576 remain with this interval; none is discarded.

Choose a minimal root \(\pi\). For \(J=L\), reduction of
\(-f(\pi Y)/\pi^{17}\) is \(Y^{17}+q(Y)\), with \(\deg q=m=20-J\),
sharing roots with every \(H_k\), \(1\le k<m\). Translate its
\(H_{m-1}\) common root to zero and scale to obtain \[
 g(Y)=Y^{17}-Y^m+\sum_{i=1}^{m-2}c_iY^i.
\] Descending solution of \(H_kg(z_k)=0\) expresses \(c_k\)
homogeneously with degree \(m-k\) in its marked witnesses. The remaining
equations \(g(z_k)=0\) have leading monomials \(z_k^{17}\), with lower
tails of total degree \(m<17\). Pairwise-coprime leading monomials give
a Gröbner basis, and the finite leading-model algebra has dimension
\(17^{m-2}\). This retains nilpotents, but not a classification of later
characteristic-zero lifts. Its translated zero need not be the original
mean.

The largest stratum \(J=L=16\) contains 124,068 markings and has
\(m=4\). Write \(g=Y^{17}-Y^4+bY^2+cY\), with \(H_2\) witness \(u\) and
\(H_1\) witness \(v\). Then \[
 b=6u^2,\qquad c=4v^3-12u^2v,
\] \[
 u^{17}+5u^4+4uv^3-12u^3v=0,\qquad
 v^{17}+3v^4-6u^2v^2=0.                                  \tag{D11}
\] For \(uv\ne0\), put \(t=v/u\), \(q(t)=-5-4t^3+12t\). Elimination
gives \(u^{13}=q(t)\) and \[
 t^{15}q(t)+3t^2-6=-4(t+8)^2(t-7)^2(t-1)^3F_5(t)F_6(t),
\] where \[
 F_5=t^5+t^4-8t^3+6t+6,\qquad
 F_6=t^6+t^4-4t^3+7t^2+8t-8
\] are irreducible over \(\mathbb F_{17}\). There are 14 distinct \(t\),
each with 13 choices of \(u\). The boundaries are one origin, 13 points
\(u=0,v^{13}=-3\), and 13 points \(v=0,u^{13}=-5\). Thus there are
\textbf{209 reduced models}, whose local lengths sum to 289, the full
algebra dimension.

Exact extension-field gcds show just one repeated root: multiplicity
four at the origin model, three when \(t=1\), and two otherwise. All
Hasse witnesses below that multiplicity lie there. Small-degree cluster
collapse makes it an exact multiple root. The original simple mean
cannot be there, and must have a simple first root. Thus every other
root in the original zero cluster has valuation exactly \(1/13\); the
original polynomial would have exactly one multiple root globally, of
multiplicity two, three, or four. Further, \[
\begin{array}{c|c}
a_2&0\text{ exactly, or valuation }2/13\\
a_{17}&0\text{ exactly, or valuation }14/13\\
a_{18}&0\text{ exactly, or valuation }15/13\\
a_{19}&\text{valuation }16/13.
\end{array}                                                    \tag{D12}
\] The linear coefficient follows from the product of the other roots.
The \(a_{17},a_{18}\) alternatives use uniqueness of the \(H_3,H_2\)
common locations: if a leading coefficient vanishes, the witness shares
the mean's simple class and equals the exact mean.

A first correction links these finite models to the original unit
coefficients. Let \(A_j\) be the unramified constants obtained by
triangular recursion at exact cube roots for unit labels and zero for
zero labels, starting with \(a_2=0,a_3=-1\). Unit-root perturbations
have value at least \(2\delta\). For a zero middle label put
\(x_j=\overline{w_j/\pi}\) and \(d_j=\overline{(a_j-A_j)/\pi}\). Then \[
 d_j=-j\bar a_{j-1}x_j\quad(\bar w_j=0),\qquad
 d_j=-\sum_{i=4}^{j-1}\binom ji d_i\bar w_j^{\,j-i}
       \quad(\bar w_j\ne0).
\] The next coefficient of the divided \(f(1)\) equation is \[
 \sum_{j=4}^{16}(C_j/17)d_j-\bar a_{16}x_{17}=0,
\] where \(x_{17}\) is the actual \(H_3\) witness. In canonical model
coordinates, write the first sum as \(\sum\lambda_jx_j\), let \(z_j\) be
the marked zero-cluster roots, and let \(\gamma\) be the original mean.
This becomes \[
 \sum\lambda_j(z_j-\gamma)+\bar a_{16}\gamma=0.
\] Every \(z_j,\gamma\) must be a root of \(g\), and \(\gamma\) is
simple. This is a finite coefficient-linked test, not an emptiness
result.

One uniform next-lift configuration is excluded: if \(a_{17}=a_{18}=0\),
not all zero-residue middle witnesses can equal the exact mean.
Otherwise set \(T=a_2\). Unit roots are congruent modulo \(17O\) to the
roots of \((X-1)(X^2+X+1+3T)\). Those roots are integral power series
over the fixed unramified quadratic base. Triangular recursion expresses
each middle coefficient modulo \(17O\) as a series \(A_j(T)\), or zero.
The divided equation gives \[
 F(T)+20a_{19}/17\in17O,\qquad
 F\in O_{\rm unr}[[T]],\quad F(0)\in17O_{\rm unr}.
\] By (D12), \(T=0\) or has value \(2/13\), while \(a_{19}/17\) has
value \(3/13\). A unit linear coefficient of \(F\) makes its value
\(2/13\); otherwise the value is at least \(4/13\), or one if \(T=0\).
Neither can cancel the value \(3/13\). This includes the vacuous case of
no zero middle labels. Other placements remain.

The \href{evidence/full/wild17/ROW8_BOUND_AND_DIVIDED_IDENTITY.md}{row-8
proof},
\href{evidence/full/wild17/blowup/COLLECTIVE_REDUCTION.md}{collective
reduction}, and
\href{evidence/full/wild17/blowup/m4/LIFT_CONSEQUENCES.md}{quartic lift
consequences} contain the full valuations and exact checks. Separate
audits cover the census, quartic model classification, and next-lift
argument. None excludes the complete \(J=L=16\) stratum or the
\(J=16,L<16\) cases.

\subsection{Cross-prime restrictions}\label{cross-prime-restrictions}

The normalization inherited from 17 need not remain integral at 2. Let
\(r\) be the total new scale, including subsequent unit normalization,
and put \(b_j=a_j/r^j\), \(t=1/r\), \(\lambda=\nu_2(t)\ge0\). Rows 4,6,7
have \(b_2=-t^2,b_{17}=0\), with \(b_3=0\) in row 4 and \(b_3=2t^3\) in
rows 6,7; their root \(t\) is simple. Row 9 has
\(b_2=b_{17}=0,b_3=-t^3\), with \(t\) exactly triple.

Row 9 cannot have 2-adic type B with \(\lambda>0\). Its four-root zero
cluster would be exhausted by simple zero and triple \(t\), so
\(F=X(X-t)^3V\), with \(V(0)\) a unit. The missing \(X^3\) coefficient
gives \[
 0=-3tV_0+3t^2V_1-t^3V_2.
\] After division by \(t\ne0\), the first term is a unit and the others
nonunits, a contradiction. This excludes a scale/type subcase, not row
9.

For rows 4,6,7, suppose type B has \(\lambda>0\) and repeated root \(u\)
in the zero cluster. Then \(F=X(X-t)(X-u)^2V\). Put
\(\mu=\nu_2(u)\ge1/2\). The missing \(X^3\) coefficient gives \[
 t(V_0-2uV_1+u^2V_2)=u(-2V_0+uV_1).
\] Initially this forces \(\lambda>1\). The zero-cluster quartic is then
\(X^4\pmod{2O[X]}\), and the type-B congruence yields
\(V\equiv X^{16}-1\pmod{2O[X]}\). Thus \(\nu(V_1)\ge1\), giving the
exact equality \(\lambda=\mu+1\). The \(G_{18}\) residue condition
forces \(\bar b_{18}=0\), while factorization gives
\(\nu(b_{18})=2\mu-1\). Hence \(\mu>1/2\), \(\lambda>3/2\), and \[
 \nu(b_2)=2\mu+2,\quad \nu(b_{18})=2\mu-1,\quad
 \nu(b_{19})=3\mu-1.
\] The \(H_2\) witness must be a unit; otherwise the size-four cluster
would collapse onto the simple mean. Returning to the original
normalization gives \(\nu_2(ru)=-1\), \(\nu_2(a_4)\in\{\infty,0,-4\}\),
and \(\nu_2(a_{19}/a_{18})=-1\). In particular, when
\(0<\lambda\le3/2\), every first-derivative common root must be in the
unit cluster.

The finite residue intersections remain nonempty in all other scale/type
cases. The coefficient-pattern/marked-assignment counts are, for rows
4,6,7, \(50/4095,58/6143,145/4608,97/3584\), respectively for A-nonunit,
A-unit, B-nonunit, B-unit. Row 9 gives \(22/1535,24/512,0/0,161/1536\).
They count the stated necessary residue system, not solutions of the
additional valuation conditions. The exact identities and their
verification are in
\href{evidence/full/cross_prime/CROSS_PRIME_RESTRICTIONS.md}{the
cross-prime proof}.

\subsection{A conditional first jet at 19 and the full finite formal
cover}\label{a-conditional-first-jet-at-19-and-the-full-finite-formal-cover}

Normalize a repeated root exactly to one at 19. The reduction is
\(X^{20}-X=X(X-1)^{19}\), with the exact mean simple. Put
\(b_j=G_j(1)\). From \(f(1)-G_{19}(1)=0\) obtain \[
 \sum_{i=2}^{19}c_i a_i=0,\qquad
 c_i=\frac1{19}\binom{19}{i-1}\in\mathbb Z.               \tag{D13}
\] The polynomial \(f(1+Y)/(1+Y)=Y^2P(Y)\) has \(P\) monic of degree 17,
with all lower coefficients divisible by 19. Every nonzero displacement
therefore has valuation at least \(1/17\). If \(b_{18}\) is a unit,
\(P(0)=190b_{18}\) has valuation one; all 17 displacements have value
exactly \(1/17\), and their normalized residues are the distinct 17th
roots of unity. Root one is exactly double, and \(a_{18}=0\), since the
\(G_{18}\) witness cannot lie in the unit cluster.

For an exact support \(I\), express coefficients by triangular equations
at witnesses \(1+\delta_j\), putting inactive coefficients zero. Let
\(\mathcal F(\delta)\) be (D13) after substitution, and
\(K_j=\partial_j\mathcal F(0)\pmod{19}\). The constant residue
condition, followed by division by a minimal displacement, gives \[
 \sum_{j\in I}K_j\zeta_j=0,\qquad
 \zeta_j\in\{0\}\cup\mu_{17},\quad\zeta_{19}=0,\qquad
\sum K_j=\bar b_{18}\ne0.                               \tag{D14}
\] For the last identity, move every formal witness by the same
parameter \(s\). Triangular homogeneity gives
\(a_i(s)=\alpha_i(1+s)^i\). Differentiating (D13) and subtracting its
constant residue equation leaves \(\bar b_{18}\). This proves the weight
sum without assuming that a selected witness actually has nonzero
displacement. A zero \(\zeta_j\) means an exact zero displacement, not
just a higher-valuation one. For example, \(I=\{12,13,15,16,17,19\}\)
has \(b_{18}=5\) and weights \((9,0,15,0,0,0)\) modulo 19. Since
\(\mu_{17}\cap\mathbb F_{19}^*=\{1\}\), the nontrivial ratio 7 prevents
nonzero phases at 12 or 15, forcing those exact witnesses to one. Other
phases remain unconstrained, and the all-zero assignment is not
prohibited. The cases \(b_{18}=0\) and higher jets remain. See
\href{evidence/full/global/P19_CLUSTER_FIRST_JET.md}{the first-jet
proof}.

Finally, every degree-20 candidate has a finite formal cover at 19. Over
\(\mathbb Z_{19}\), take variables \(a_2,\ldots,a_{19}\) and
\(w_2,\ldots,w_{18}\), put \(w_{19}=1\), and impose all 36 equations
\(f(w_j)=G_j(w_j)=0\), \(2\le j\le19\). The special fibre has
\(a_{19}=-1\), hence \(f=X(X-1)^{19}\). Every witness is zero or one,
and triangular equations determine all coefficients. Conversely every
word of 17 bits works. There are exactly \(2^{17}\) marked geometric
points, all rational over \(\mathbb F_{19}\).

The special fibre is finite-type and zero-dimensional, hence Artinian
with nilpotents retained. Complete the local ring at a point, obtaining
\(R\). Since \(R/19R\) is Artinian, a power of the maximal ideal lies in
\(19R\), so its topology agrees with the 19-adic topology. Lift a basis
of \(R/19R\). Successive approximation and completeness express every
element as a \(\mathbb Z_{19}\)-linear combination of those lifts. Thus
\(R\) is finite over \(\mathbb Z_{19}\), with no flatness or reducedness
assumption.

Nonintegral normalized candidates cannot escape this cover. For a point
already satisfying \(w_{19}=1\), scale by a minimum-valuation root
\(\eta\). Its scaled reduction is \(X^{20}+\bar a_{19}X\), with
\(\bar a_{19}\ne0\). The chosen repeated root \(1/\eta\) cannot reduce
to zero, since the scaled \(G_{19}\) does not vanish there. Therefore
\(\eta\) is a unit and the original point was integral. Every algebraic
candidate maps to one completed local ring, including arbitrary ramified
extensions. Their generic fibres are finite algebras, and the whole
normalized generic fibre is consequently finite as well.

No local lengths, field factors, numerical ramification bounds, or
candidate list have been computed for this cover. Finiteness does not
show that any generic fibre vanishes. This standard argument includes
all 576 exceptional row-8 markings and every other branch; it does not
complete their lifting computations. See the
\href{evidence/full/wild17/blowup/COLLECTIVE_REDUCTION.md}{formal-cover
proof} and
\href{evidence/full/wild17/blowup/COLLECTIVE_REDUCTION_AUDIT.md}{independent
audit}.

The conclusion remains limited: row 3 is excluded completely; the other
eight characteristic-17 branches have the restrictions and partial
exclusions stated here. The complete support calculations, finite
leading algebras and generic norm criterion give exact targets for
further work. None proves unrestricted degree 20, and a degree-20 result
alone would not establish the all-degree conjecture.

\section{All-degree deductions, countermodels, and unresolved
implications}\label{all-degree-deductions-countermodels-and-unresolved-implications}

This section records the completed all-degree investigations. It
contains proved necessary conditions, exact reformulations, and
countermodels to specified weakened arguments. It does not prove the
Casas--Alvero conjecture, establish novelty, or assess tractability.
Positive-characteristic examples below are not characteristic-zero
counterexamples. Examples satisfying all but one common-root condition
are explicitly identified as such.

In characteristic zero, write a centered monic polynomial as

\[
 f(X)=\sum_{i=0}^n\binom ni a_iX^{n-i},\qquad a_0=1,
 \qquad G_j(X)=\sum_{i=0}^j\binom ji a_iX^{j-i}.
\]

Thus \(G_j=f^{(n-j)}/(n!/j!)\) and \(G'_j=jG_{j-1}\). Centering means
\(a_1=0\). For a CA polynomial the mean is a root, so \(a_n=0\) as well.
The CA conditions are that \(f\) and each \(G_j\), \(1\le j<n\), have a
common root. In arbitrary characteristic, \(H_j\) denotes the Hasse
derivative, defined by \(f(X+T)=\sum_jH_jf(X)T^j\).

\subsection{What a characteristic-zero contraction repair must
prove}\label{what-a-characteristic-zero-contraction-repair-must-prove}

The first investigation concerns the elimination step associated with
\href{https://arxiv.org/html/2501.09272v2}{Ghosh's v2 proof claim,
Sections 4.1--4.2}. The algebraic statements here do not assume
Proposition 3.3. The application to that paper retains its separate
determinantal-height input; this section does not reprove the prequel
supplying that input.

\textbf{Proposition 1 (contraction equivalence).} Let
\(R=k[x_1,\ldots,x_m]\), \(S=R[t]\), and

\[
 I=(tf_i+g_i:1\le i\le m),\quad J=(f_1,\ldots,f_m),
 \quad D=I_2\begin{pmatrix}f_1&\cdots&f_m\\g_1&\cdots&g_m\end{pmatrix}.
\]

Assume that: (i) some \(L=t+\ell\), \(\ell\in R\), satisfies
\(\operatorname{ht}(I+(L))=m+1\); (ii) all minimal primes of \(D\) have
height \(m-1\), and \(R/D\) has no embedded associated primes; (iii)
\(J\) is proper and has height at least \(m-1\). Then

\[
             I\cap R=D\quad\Longleftrightarrow\quad
             \operatorname{ht}J=m.                         \tag{1}
\]

\textbf{Proof.} The minors belong to \(I\cap R\). If
\(\operatorname{ht}J=m\), each minimal prime \(\mathfrak p\) of \(D\)
omits some \(f_l\). After localizing there, eliminate \(t\) with
\(tf_l+g_l\). The resulting equations are exactly the minors involving
column \(l\), which generate all minors because \(f_l\) is a unit. Thus
contraction agrees with \(D\) at every minimal prime. A nonzero
submodule of \(R/D\) has an associated prime among those of \(R/D\); the
absence of embedded primes therefore makes the kernel of
\(R/D\longrightarrow S/I\) zero.

Conversely, suppose \(\operatorname{ht}J=m-1\), and choose a minimal
prime \(\mathfrak p\) of \(J\) of that height. Since \(D\subset J\), it
is also minimal over \(D\). Some \(g_l\) is a unit in
\(R_{\mathfrak p}\): otherwise \(I+(L)\subset\mathfrak pS+(L)\),
contradicting its required height. In the Artinian ring
\(A=R_{\mathfrak p}/D_{\mathfrak p}\), every \(f_i\) is nilpotent.
Consequently \(tf_l+g_l\) is a unit of \(A[t]\), by a finite
geometric-series inverse. Thus \(IS_{\mathfrak p}=S_{\mathfrak p}\),
whereas \(D_{\mathfrak p}\) is proper. Contraction fails at this
component. These are the only possible heights of proper \(J\).
\(\blacksquare\)

In the paper's elementary-symmetric family, the right side of (1) is its
lower-degree CA regular-sequence condition; the upper-degree CA
assumption provides the linear completion. Accordingly, proving this
contraction identity uniformly would prove the missing descent, rather
than replace it with a weaker routine lemma. A component on which every
\(f_i\) vanishes and some \(g_i\) is invertible is precisely the
component lost by the finite-\(t\) elimination chart.

\textbf{Countermodel 1 (generic reducedness is unavailable).} In the
elementary-symmetric family on \(x,y,z\), put

\[
 a=xyz,\ b=xy+xz+yz,\ c=x+y+z,\qquad
 f=(a,b,c),\quad g=(0,a,b).
\]

The determinantal ideal is \(D=(a^2,ab,b^2-ca)\), with radical
\((xy,xz,yz)\). At \(\mathfrak p=(x,y)\), \(z\) is a unit. The only
order-two initial form among the displayed generators is
\(z^2(x^2+xy+y^2)\); the others start in orders three and four. The
order-two form \(zxy\) of \(a\) is not its scalar multiple over
\(\mathbb Q(z)\). Hence \(a\notin D_{\mathfrak p}\), although
\(a^2\in D_{\mathfrak p}\). The quotient is not generically reduced.
This is not a contraction counterexample: \(c\) is a unit at this prime,
and exact elimination confirms contraction in this example.

\textbf{Countermodel 2 (the general structural hypotheses do not repair
it).} Over \(R=\mathbb Q[u,v,z,s]\), take

\[
 f=(u^4,uvs,v^2,z),\qquad g=(-u^5,s^4,2v^3,3z^2).
\]

These have the descending degree pattern and the first-column relation
of the proposed argument. At \(\mathfrak p=(u,v,z)\),
\(J_{\mathfrak p}=(u^4,uv,v^2,z)\) still needs four generators, equal to
the global number. Because \(g_2=s^4\) is a unit, the column-2 minors
give

\[
 D_{\mathfrak p}=(u^4,v^2,z),\qquad
 R_{\mathfrak p}/D_{\mathfrak p}
 =\mathbb Q(s)[u,v]_{(u,v)}/(u^4,v^2).
\]

Here \(uv\ne0\) is square-zero and \(s^4+uvst\) is a polynomial unit.
Contraction therefore fails despite the local/global generator equality.
The upper equations form a regular sequence: modulo \(t\), their
constants generate \((u^5,s^4,v^3,z^2)\), of height four. The
determinantal ideal has expected height three: at primes containing
\(J\) this follows from \(\sqrt J=(u,v,z)\); elsewhere a unit \(f_i\)
eliminates \(t\) from the upper complete intersection. Thus the standard
expected-height determinantal criterion also gives Cohen--Macaulayness.

Even all six linear completions \(L_r=u+v+z+s+t-6r\),
\(r\in\{u,v,z,s,t,0\}\), have height five. For a direct check, when
\(u,v\ne0\), normalize \(u=t=1\), obtaining \(v=-1/2\),
\(z\in\{0,-1/3\}\), and \(s=0\) or \(s^3=1/2\). None of the six linear
equations vanishes; each prescribes an incompatible rational value of
\(s\). When either \(u\) or \(v\) is zero, \(s=0\) and the remaining
factors give the same conclusion. Exact ideal-dimension checks accompany
this case analysis. This is an abstract algebraic model, not a CA
configuration. It shows why derivative-specific identities are essential
to any repair.

Evidence: \path{evidence/full/global/REPAIR_EQUIVALENCE.md},
\texttt{actual\_nilpotents.sing}, \path{descending_degree_model.sing},
and their \texttt{*-verification.txt} outputs in that directory. The
unresolved assertion is the exclusion of the bad components in the
special derivative-linked family.

\subsection{Cluster counts and an occupancy countermodel at every
scale}\label{cluster-counts-and-an-occupancy-countermodel-at-every-scale}

\textbf{Proposition 2 (derivative counts in a separated cluster).} Work
over an algebraically closed nonarchimedean valued field. A separated
closed disk contains \(m\) roots of \(f\), counted with multiplicity.
Rescale it to the unit disk and divide by a nonzero scalar, obtaining
\(Q(Y)=A(Y)B(Y)\), where \(A\) is monic of degree \(m\), all its roots
are integral, and \(\overline B=1\). Then \(H_mf\) has no root in the
original disk. If \(H_j\overline A\ne0\) has degree \(d\), \(H_jf\) has
exactly \(d\) roots there. In particular the count is \(m-j\) when
\(j<m\) and \(\binom mj\) is a residue-field unit.

\textbf{Proof.} Reduction commutes with Hasse differentiation, so
\(\overline{H_mQ}=1\), proving nonvanishing on integral arguments. An
integral polynomial of Gauss norm one whose reduction has degree \(d\)
has exactly \(d\) integral roots: factor its integral-root factors as
\(Y-\beta\) and its nonintegral-root factors as \(1-Y/\beta\). The
latter reduce to 1; the remaining scalar is a unit. Reduction therefore
counts the integral roots exactly. Apply this to \(H_jQ\) and undo the
affine change. If the entire reduction is zero, this argument makes no
root-count claim. \(\blacksquare\)

These counts alone say nothing about which derivative roots are also
roots of \(f\). The following exact family prevents a general
occupancy-only induction from supplying that missing information.

\textbf{Proposition 3 (digit-block Hasse--CA models).} In characteristic
\(p\), choose distinct nonnegative integers \(e_i\), digits
\(1\le d_i<p\), and distinct root positions \(a_i\). Put

\[
 N=\sum_i d_ip^{e_i},\qquad P(X)=\prod_i(X-a_i)^{d_ip^{e_i}}.
\]

For \(j=\sum_ej_ep^e\), with unused degree digits zero,

\[
 H_jP=\prod_i\binom{d_i}{j_{e_i}}
            (X-a_i)^{(d_i-j_{e_i})p^{e_i}}                 \tag{2}
\]

when every \(j_e\le d_e\), and is zero otherwise. Consequently \(P\) is
Hasse--CA for arbitrary positions, and is nontrivial when at least two
positions occur.

\textbf{Proof.} Expand each Taylor factor using Frobenius:
\((X-a_i+T)^{d_ip^{e_i}}
=\sum_{k=0}^{d_i}\binom{d_i}k(X-a_i)^{(d_i-k)p^{e_i}}T^{kp^{e_i}}\). The
possible exponents have unique carry-free base-\(p\) expansions, proving
(2). For \(0<j<N\), either the derivative is zero or a root factor
remains. \(\blacksquare\)

For any nonempty proper subset \(S\) of positions, let
\(m_S=\sum_{i\in S}d_ip^{e_i}\). Every derivative order \(1\le j<m_S\)
has a common root among those positions: removing all their factors
requires \(j\ge m_S\). At order \(m_S\) all and only their factors are
removed, so that derivative has no root there. Thus every separated
cluster has the maximum possible initial occupancy and satisfies
Proposition 2's counts. Nevertheless, every cluster containing distinct
positions has a nontrivial local Hasse--CA polynomial of its own degree
by the same construction.

There can be arbitrarily many such nested clusters: over
\(\overline{\mathbb F_p((t))}\), use positions \(1,t,\ldots,t^r,0\) with
distinct digit-position multiplicities \(p^i\). Repeated rescaling does
not produce a good local degree. In degree 20, two-position examples
occur at primes \(2,3,7,11,13,17,19\), using the digit decompositions
\(4+16,2+18,6+14,9+11,7+13,3+17,1+19\). The last example is
\(X(X-1)^{19}=X^{20}-X\) in characteristic 19.

This does not obstruct every characteristic-zero cluster proof. Indeed
\(X(X-1)^{19}\) over \(\mathbb Q\) satisfies orders 1 through 18 but
fails order 19, where the derivative is \(20X-19\). Divisibility hidden
by reduction supplies the missing obstruction. A valid use of cluster
collapse must establish both the required exact witness occupancy and a
good local residue degree, or add a lifting argument.

There is also a finite-prime limitation. For a fixed algebraic
polynomial, put all distinct roots in one number field. Outside finitely
many primes, all nonzero root differences are units. Every disk
containing two distinct positions then contains every position. Only
finitely many primes can therefore supply a proper separated cluster
with distinct roots. This follows from the finite prime support of each
difference and its denominator; it does not rule out success at an
exceptional prime.

Evidence: \path{evidence/full/cluster_tree/COUNTERMODEL_AND_COUNTS.md}
and \texttt{verify\_digit\_blocks.py}; the conditional collapse lemma is
in \path{evidence/full/two_adic/CLUSTER_COLLAPSE.md}. The replay checks
direct Taylor expansions, every derivative, and subset occupancies for
the displayed models.

\subsection{Integral reconstruction and a linked Newton
bound}\label{integral-reconstruction-and-a-linked-newton-bound}

The exact integral recursion is

\[
       G_j(z)=j\int_{r_j}^{z}G_{j-1}(t)\,dt,
       \qquad G_j(r_j)=0.                                \tag{3}
\]

Polynomial integrals are path independent. The additional equations
\(f(r_j)=0\) are indispensable: convex-hull membership is insufficient
even when all but one incidence is exact.

\textbf{Countermodel 3.} For every \(n\ge4\),

\[
 f_n=X(X-1)^{n-2}(X+n-2),\qquad
 G_j=(X-1)^{j-2}\left(X^2+(j-2)X-\frac{(j-1)(n-j)}n\right)
                                                               \tag{4}
\]

for \(2\le j\le n\), and \(G_1=X\). These identities follow by
differentiation and the terminal equality \(G_n=f_n\). The polynomial is
centered and real-rooted, with simple mean root 0 and repeated root 1.
It satisfies every CA condition except that for \(G_2=X^2-(n-2)/n\).
Choose nodes \(r_1=0\), \(r_2=\sqrt{(n-2)/n}\), and \(r_j=1\) for
\(j\ge3\), including the terminal node. All lie in \([0,1]\), all but
one are final roots, and they satisfy (3). Yet the final root radius is
\(n-2\), while the node radius is 1. These hypotheses imply neither
contraction nor a degree-independent bound by a constant times the node
radius. The example does not refute an estimate using all incidences.

\textbf{Proposition 4 (first nonzero coefficient witness).} For a
nontrivial centered CA polynomial, let \(d\) be least with \(a_d\ne0\),
choose a common root \(r\) of \(G_d=X^d+a_d\) and \(f\), and let
\(m_0,m_r\) be the multiplicities of 0 and \(r\). Then \(2\le d<n\),
\(r\ne0\), and, writing \(R\) for the final root radius,

\[
 \left(d\binom nd-m_r\right)|r|^d
       \le(n-m_0-m_r)R^d.                                \tag{5}
\]

\textbf{Proof.} Newton's identity and \(a_d=-r^d\) give
\(\sum_\alpha\alpha^d=d\binom nd r^d\), with multiplicities. Removing
the zero roots and the \(m_r\) copies of the actual root \(r\), and
applying the triangle inequality, proves (5). Its left coefficient
exceeds its right coefficient, so \(0<|r|<R\). Equality requires every
remaining root to have modulus \(R\) and its \(d\)th power to have the
argument of \(r^d\). \(\blacksquare\)

In particular all nonzero roots cannot lie on one circle about the mean.
The estimate does not iterate by translating to \(r\): the new first
coefficient occurs at index 1, and its linear derivative selects the old
mean. No decreasing quantity across these changes of center has been
proved.

For the special normalization \(a_2=-1,a_3=2\) with simple roots 0 and
1, Newton's identities also give

\[
 R^2\ge\frac{n(n-1)-1}{n-2},\qquad
 R^3\ge\frac{n(n-1)(n-2)+1}{n-2}.                      \tag{6}
\]

The degree-five polynomial \(X(X-1)(X^3+X^2-9X+11)=X^5-10X^3+20X^2-11X\)
has these normalized \(G_2,G_3\), and simple roots 0 and 1, but has a
nonreal pair and is squarefree. It satisfies the other three derivative
conditions and fails the first. Thus real-rootedness of \(G_3\) cannot
be propagated backward. If the original roots were collinear, a double
root of \(G_3\) would instead force multiplicity at least \(n-1\) in
\(f\): for a real-rooted polynomial every non-inherited derivative root
is simple, as follows from \((P'/P)'=-\sum m_i/(x-r_i)^2<0\). Iterating
and rotating proves this restricted statement, not an exclusion of
complex configurations.

Evidence: \path{evidence/full/global_analytic/INTEGRAL_OBSTRUCTION.md}
and \path{verify_integral_countermodels.py}. The arbitrary-degree proofs
above are supplemented by exact replays of (4) for degrees 4--40 and of
the degree-five example. The unresolved step is an additional
all-incidence estimate giving an incompatible upper bound or an actual
descent.

\subsection{What the product formula
yields}\label{what-the-product-formula-yields}

Now assume the coefficients and roots are algebraic, in a common number
field \(K\). For each place use the usual absolute value and weights
\(w_v=[K_v:\mathbb Q_v]/[K:\mathbb Q]\), without separately squaring
complex moduli. Put

\[
 R_v=\max_i|\alpha_i|_v,\qquad H_R=\prod_vR_v^{w_v}.
\]

This root-vector projective height is unchanged by one global root
scaling. Fix the same algebraic witness \(r\) from Proposition 4 at all
places and write \(B=d\binom nd-m_r\), \(M=n-m_0-m_r\).

\textbf{Proposition 5 (height lower bound).} One has \(M>0\) and

\[
 H_R\ge(B/M)^{1/d}>1,
 \qquad H_R\ge
 \left(\frac{n(n-1)-1}{n-2}\right)^{1/(n-1)}.             \tag{7}
\]

\textbf{Proof.} The exact identity underlying (5) is
\(Br^d=\sum_{\alpha_i\ne0,r}\alpha_i^d\); if \(M=0\) it is impossible.
At every infinite place it gives \(R_v/|r|_v\ge(B/M)^{1/d}\). At finite
places the stronger relevant lower bound is simply \(R_v/|r|_v\ge1\),
because \(r\) is an actual root. Multiply and use the product formula on
\(r\), and the sum 1 of the infinite-place weights. For the second
estimate, \(d\binom nd\ge n(n-1)\); the ratio is minimized at
\(m_0=m_r=1\), and \(d\le n-1\). \(\blacksquare\)

The product formula does not set \(H_R=1\). Different roots can maximize
the radius at different places. Normalization by a different root at
each place is not multiplication by one global element.

\textbf{Proposition 6 (coefficient-height comparison).} Define
\(Q_v=\max_{1\le i\le n}|a_i|_v^{1/i}\) and \(H_A=\prod_vQ_v^{w_v}\).
Then \(Q_v=R_v\) at finite places, and

\[
 H_A\le\kappa H_R,\qquad
 \kappa=\sqrt{\frac{(n-m_0)(n-m_0-1)}{n(n-1)}}<1.         \tag{8}
\]

\textbf{Proof.} A common witness \(\beta_j\) gives
\(a_j=-\sum_{i<j}\binom ji a_i\beta_j^{j-i}\). Nonarchimedean induction
therefore yields \(|a_j|_v\le R_v^j\). The monic root-radius formula
gives \(R_v=\max_i|\binom ni a_i|_v^{1/i}\le Q_v\), proving equality. At
infinite places the elementary symmetric functions, after removing
\(m_0\) zero roots, give \(|a_i|_v\le[\binom{n-m_0}i/\binom ni]R_v^i\).
The successive factors of this ratio decrease. Since \(a_1=0\), their
geometric means for \(i\ge2\) are at most \(\kappa\). Multiplication
proves (8). \(\blacksquare\)

These are comparisons between different heights, not a self-decreasing
height. Scaling to \(r=1\) gives \(a_d=-1\), so \(H_A\ge1\); it merely
requires a larger root height. The root-difference height
\(\prod_v(\max_{i,j}|\alpha_i-\alpha_j|_v)^{w_v}\) is invariant under
translation and global scaling, so translating to a new root cannot
strictly decrease that alternative either.

\textbf{Countermodel 4.} The exact polynomial

\[
 F=X^6-15X^4+20X^3-6X
   =X(X-1)^2(X^3+2X^2-12X-6)
\]

has normalized coefficients \((1,0,-1,1,0,-1,0)\). It satisfies four of
five CA conditions, including the first nonzero-coefficient witness and
ordinary derivative, but \(\gcd(F,G_3)=1\). At every finite place its
roots are integral by monicity, and the exact root 1 gives
\(R_v=Q_v=1\). The cubic has roots in \((-5,-4),(-1,0),(2,3)\), by
endpoint signs, so \(H_R\in(4,5)\) while \(H_A=1\). Bound (7) here is
\(H_R\ge\sqrt{28/3}\), which is compatible with the model.

Moreover \(F_t=X(X-1)^2(X^3+2X^2-12X-6+60t)\), \(t\in\mathbb Z\), has
normalized coefficients \((1,0,-1,1+3t,-8t,-1+10t,0)\). It retains
witnesses for \(G_1,G_2,G_5\), a simple mean root, a double root 1, and
all finite-place coefficient bounds, while
\(H_R(F_t)\ge|60t-6|^{1/3}\to\infty\). The remaining incidences are not
asserted. Thus the stated weakened data supply no height upper bound.

Evidence: \path{evidence/full/global_height/ALL_PLACE_BOUND.md} and
\texttt{verify\_height\_models.py}. Integrality at all finite places
follows from monicity and the exact unit root, not from sampling primes.
An all-incidence height inequality beyond (7)--(8) remains missing.

\subsection{Exact matrix flags and limits of spectral
shortcuts}\label{exact-matrix-flags-and-limits-of-spectral-shortcuts}

For \(D=\operatorname{diag}(\alpha_1,\ldots,\alpha_n)\), there is an
exact nested orthogonal compression flag with characteristic polynomials
\(G_1,\ldots,G_n\). The input is Pereira's theorem: every complex matrix
has a unit trace vector \(v\), satisfying
\(v^*A^kv=\operatorname{tr}(A^k)/m\) for \(0\le k<m\), and compression
to \(v^\perp\) has polynomial \(p_A'/m\). This includes nonnormal
matrices. The primary statement is explicitly recalled in
\href{https://www.numdam.org/item/CRMATH_2005__341_11_651_0.pdf}{Pereira
(2005), Definitions 3--4 and Proposition 5}.

For completeness, the identity follows from
\(p_{A|v^\perp}(z)=p_A(z)v^*(zI-A)^{-1}v\): the trace-vector moments and
Cayley--Hamilton identify the resolvent expression with
\(p_A'(z)/(mp_A(z))\). Iterating the existence theorem gives the flag;
compression along nested orthogonal subspaces is associative. In an
adapted basis the full matrix is normal, every leading block has
polynomial \(G_j\), and centering makes every diagonal entry zero. CA is
exactly the shared- spectrum condition between every proper block and
the full matrix. This is a reformulation, not an exclusion.

\textbf{Proposition 7.} The first compression along the flat vector
\(v=(1,\ldots,1)/\sqrt n\) is normal if and only if the roots lie on an
affine line.

\textbf{Proof.} Write \(D\) in the decomposition
\(v^\perp\oplus\mathbb Cv\) as
\(\begin{pmatrix}B&b\\c^*&\mu\end{pmatrix}\), where
\(\mu=\operatorname{tr}D/n\), \(b=(D-\mu)v\), and
\(c=(D^*-\overline\mu)v\). Normality gives \(BB^*-B^*B=cc^*-bb^*\). If
\(b=0\), all roots coincide. Otherwise equality of these rank-one
matrices is equivalent to \(c=\eta b\), \(|\eta|=1\), or
\(\overline{\alpha_i-\mu}=\eta(\alpha_i-\mu)\) for every \(i\). This is
exactly collinearity. \(\blacksquare\)

\textbf{Proposition 8.} For \(n\ge5\), no fixed \(n\times2\) isometry
\(V\) represents \(G_2\) as the characteristic polynomial of \(V^*DV\)
for all root tuples.

\textbf{Proof.} Cauchy--Binet and comparison of each monomial
\(\alpha_i\alpha_j\) force every squared two-row minor of \(V\) to be
\(1/\binom n2\). Summing minors gives each row squared norm \(2/n\).
After row normalization, the rank-one Hermitian projectors have pairwise
Hilbert--Schmidt products \(t=(n-2)/(2(n-1))\) off the diagonal and 1 on
it. Their Gram matrix \((1-t)I+tJ\) is positive definite, so the \(n\)
projectors are real-linearly independent. The ambient Hermitian
\(2\times2\) space has dimension four, a contradiction. \(\blacksquare\)

This rules out a root-independent universal flag, not the root-dependent
flag just constructed. Shared spectrum also does not imply that
compressed eigenvectors come from original eigenspaces, even in an exact
derivative flag.

\textbf{Countermodel 5.} Start with the normal operator
\(0\oplus C_{n-1}\), where \(C_{n-1}\) cyclically permutes
\(u_1,\ldots,u_{n-1}\), and the separate zero vector is \(e_0\). Put
\(b=1/\sqrt n\), \(a=\sqrt{(n-1)/n}\), \(w=ae_0+bu_{n-1}\),
\(v=-be_0+au_{n-1}\). In the ordered basis \(u_1,\ldots,u_{n-2},w,v\),
the characteristic polynomials of its leading blocks are

\[
 p_n=z^n-z,\qquad p_{n-1}=z^{n-1}-1/n,\qquad
 p_j=z^j\quad(1\le j\le n-2).                          \tag{9}
\]

Indeed the penultimate block is a weighted cycle with edge-product
\(b^2=1/n\), and the smaller blocks are nilpotent shifts. These are
exactly the normalized derivatives of \(z^n-z\). The original eigenvalue
0 is simple, with eigenvector \(e_0=aw-bv\), orthogonal to every one of
the first \(n-2\) subspaces. Yet their blocks have 0 as their sole
eigenvalue, with algebraic multiplicity \(j\). The final required
incidence fails: the penultimate roots satisfy \(z^{n-1}=1/n\), whereas
the original nonzero roots satisfy \(z^{n-1}=1\). Thus this is another
all-but-one-incidence countermodel, not a CA polynomial.

Evidence: \path{evidence/full/global_matrix/COMPRESSION_AUDIT.md} and
\texttt{verify\_compressions.py}. The replay checks the explicit
dimension-four matrix over \(\mathbb Q(\sqrt3)\), its normality,
intermediate nonnormality, all characteristic polynomials, trace-vector
moments, and kernel vector. The arbitrary-degree arguments are algebraic
proofs. A matrix proof would still need an obstruction coupling all
common-spectrum conditions with all trace-vector equations; isolated
eigenvalue sharing does not supply it.

\subsection{Review boundary and replay
map}\label{review-boundary-and-replay-map}

The four Python replay directories listed above retain
\texttt{verification.json} and \path{verification-optimized.json}; they
report \texttt{PASS} under ordinary and optimized execution. The
algebraic repair dossier retains exact Singular outputs. These checks
certify the displayed finite calculations, while the arguments in this
section carry the arbitrary-degree quantifiers. They do not certify the
conjecture, the unproved contraction descent, full cluster occupancy, a
height descent, or an all-incidence spectral exclusion.

The degree-20 prime-19 first-jet calculation is separate:
\path{evidence/full/global/P19_CLUSTER_FIRST_JET.md} and
\texttt{verify\_p19\_first\_jet.py}. Its valuation and phase
restrictions are not all-degree theorems and are not needed for the
statements above.

\section{Remaining obligations and review
priorities}\label{remaining-obligations-and-review-priorities}

The unrestricted degree-20 problem still requires an exclusion of every
candidate in the eight surviving characteristic-17 branches. The branch
tables describe ordinary reductions of polynomials; many normalized
coefficients are invisible there. All admissible placements of
derivative witnesses must also be included. Finite-field nonexistence in
a chosen specialization, or a count of local algebra dimensions, cannot
substitute for that coverage.

The tame branches now have explicit lifting descriptions, and the
completed support exclusions make part of their finite search decisive.
Their largest systems dominate the remaining cost. Counts of excluded
supports do not measure a proportional distance to a full proof. In the
wild branches, leading models and valuation constraints still leave
genuine lifting questions. No deduction in this report eliminates all
eight branches.

For the all-degree conjecture, the missing requirement is a uniform
argument. The contraction repair is not such an argument because its
decisive injectivity statement retains the lower-degree Casas--Alvero
problem. The moment inequalities provide lower bounds without a
contradictory upper bound. The height argument compares two different
projective heights. The matrix flag represents the exact derivative
chain, but the countermodel prevents an inference of inherited
eigenvectors from shared eigenvalues. None of these failures rules out a
stronger future argument; none supplies that argument now.

\subsection{Work completed for this
review}\label{work-completed-for-this-review}

The closeout consolidates the existing proofs and independently checked
certificates, refreshes the primary-source attributions, and makes a
final bounded attempt on the next complete row-9 batch. The
computation's actual outcome and limits are recorded with the
case-exclusion results. All supplied replay commands have explicit
scope; the archive preserves unfinished and negative results as such. No
additional open-ended proof campaign is a dependency of this review
package.

\subsection{Questions for mathematical
review}\label{questions-for-mathematical-review}

The most consequential review tasks are the following.

\begin{itemize}
\tightlist
\item
  Check the normalization, coefficient integrality, simple-mean input,
  and transfer from algebraic incidence points to the local fields used
  in each exclusion. In particular, test every change from a residue
  statement to an identity in the valuation ring.
\item
  Check the completeness of the nine-seed classification and the
  subsequent marked covers. For row 4, examine the signed quadratic
  charts and the coefficientwise divided obstruction. For rows 6, 7 and
  9, examine the Jacobian and uniqueness arguments over ramified ambient
  extensions.
\item
  Check the row-8 scale bounds, retention of the exceptional
  assignments, treatment of nilpotents in the leading algebra, and
  identification of the original mean after translation. The uniform
  next-lift exclusion depends on exact witness locations, not just equal
  residues.
\item
  Check the sparse theorem's complete support and branch coverage. Its
  short valuation proof, rather than the optional large integer
  resultant, is the main proof of the final six-term family exclusion.
\item
  Independently assess novelty and significance. The unavailable-source
  overlap recorded in the seven-term dossier remains material. A
  successful replay, an explicit certificate, and an internal audit do
  not establish either historical priority or sufficient importance for
  publication.
\end{itemize}

\texttt{REVIEW\_GUIDE.md} gives a shorter reading order and the
principal claim-to-file map. The complete original proof notes remain
available for a reviewer who wants to examine each intermediate lemma
and experiment.

\subsection{Provenance and
availability}\label{provenance-and-availability}

This is an unpublished research report prepared with AI assistance. The
arguments and computations were developed in a Codex task using
additional AI agents for bounded derivations and adversarial checks.
Those agents are not independent human referees. Python, exact integer
and finite-field arithmetic, Singular, and selected symbolic libraries
were used as identified in the evidence files. No empirical data, human
participants, or animal subjects are involved.

The package contains the report source, proof notes, scripts, finite
certificates, recorded outputs, and a file-integrity manifest.
Referenced third-party manuscripts are linked rather than redistributed.
Authorship, funding, and conflict-of-interest declarations for any
future submission have not been supplied and are not invented here. No
submission, publication, author contact, or external release was
performed as part of this closeout.

\section{References, versions, and source
limits}\label{references-versions-and-source-limits}

Checked 23 September 2026. These entries distinguish a primary source's
claim from a theorem independently checked in this campaign. ``Current''
means the record retrieved on this date; it does not assert that no
other manuscript, public comment, or private calculation exists. Version
histories were refreshed for the arXiv entries below during preparation
of this review. The thesis and supporting full-text passages were
inspected in earlier bounded rounds on the same date, as recorded in the
source notes.

\subsection{Primary mathematical
sources}\label{primary-mathematical-sources}

\textbf{{[}Ghosh-proof{]} Soham Ghosh, \emph{Proof of the Casas-Alvero
conjecture}.}
\href{https://arxiv.org/abs/2501.09272v2}{arXiv:2501.09272v2},
\href{https://arxiv.org/html/2501.09272v2}{version-pinned full text}.
First submission 16 January 2025; v2, 21 March 2026, 01:41:44 UTC. The
current record shows v2 and the comment ``Major revisions''; it shows no
later revision, withdrawal, or journal reference. The captured PDF's
internal date is 24 August 2026. The campaign audit addresses
Proposition 3.3, its conormal cancellation, and the polynomial-unit
inference in Lemma 4.5. The abstract's full-conjecture claim is not
accepted as a result established by this review. Source hashes and the
byte-identity record are retained in
\href{evidence/audit/SOURCE_IDENTITY.json}{SOURCE\_IDENTITY.json}.

\textbf{{[}Ghosh-finiteness{]} Soham Ghosh, \emph{A finiteness result
towards the Casas-Alvero Conjecture}.}
\href{https://arxiv.org/abs/2402.18717v3}{arXiv:2402.18717v3},
\href{https://arxiv.org/html/2402.18717v3}{full text}. First submission
28 February 2024; v3, 14 January 2025. The
\href{https://www.press.jhu.edu/journals/american-journal-mathematics}{American
Journal of Mathematics publisher's accepted-papers list} records this
title as accepted 7 April 2026. This is the earlier finiteness paper,
not the claimed full proof. The latter uses its height/finiteness
inputs. Acceptance is bibliographic evidence, not an independent
verification of every imported lemma. This review does not claim a
complete audit of the prequel.

\textbf{{[}CLO{]} Wouter Castryck, Robert Laterveer, Myriam Ounaïes,
\emph{Constraints on counterexamples to the Casas-Alvero conjecture, and
a verification in degree 12}.}
\href{https://arxiv.org/abs/1208.5404v1}{arXiv:1208.5404v1}, 27 August
2012; \href{https://arxiv.org/html/1208.5404}{full text}. Relevant
locations are Theorem 2 (degree \(p+1\), mean-root and determinant
restrictions), Theorem 13 (distinct-root bound), Proposition 15
(prime-power derivative placement), and the scenario/descendant
treatment in the computational sections. These are established
predecessor methods, not campaign contributions. Their demonstrated
full-degree computation is degree 12. In degree 20, derivative orders
\(1,19\), orders \(4,16\), and the simultaneous triple \(5,10,15\)
supply forbidden shared-root configurations. The appropriate hypotheses
and normalization must be retained when translating these into
coefficient conditions.

\textbf{{[}deFrutos{]} Rosa María de Frutos Marín, \emph{Perspectivas
aritméticas para la Conjetura de Casas-Alvero}.} Doctoral thesis,
Universidad de Valladolid, repository year 2013; advisor Antonio
Campillo López.
\href{https://uvadoc.uva.es/bitstream/10324/3602/1/tesis367-130927.pdf}{Primary
PDF},
\href{https://uvadoc.uva.es/handle/10324/3602?show=full}{repository
metadata}, DOI
\href{https://doi.org/10.35376/10324/3602}{10.35376/10324/3602}. The
title-page date is 21 December 2012, the university catalog defense date
is 11 June 2013, and repository availability is 27 September 2013; these
describe distinct events. Sections 2.3 and 3.5, especially Proposition
3.5.5, contain the older sparse/two-visible-coefficient criterion used
in the campaign's support sieve. Indexed primary PDF text was retrieved
and inspected; local byte download and screenshot retrieval were
unsuccessful, so no local PDF hash or visual-page check is claimed. A
short companion is her official 2015 abstract,
\href{https://www.singacom.uva.es/JTN2015/contribuciones/ordinarias/frutos.pdf}{\emph{Un
problema sobre números combinatorios}}, at the Valladolid number-theory
meeting, 29 June--3 July 2015.

\textbf{{[}Massri{]} Cesar Massri, \emph{The Casas-Alvero conjecture for
three recycled roots in degree 20}.}
\href{https://arxiv.org/abs/1806.09561v6}{arXiv:1806.09561v6},
\href{https://arxiv.org/html/1806.09561v6}{full text}. First submission
25 June 2018; v6, 25 August 2023. Some earlier versions are marked
withdrawn; the current v6 title and scope should be cited, not an
earlier full-proof claim. Relevant passages are Remark 7.4, the proof of
Theorem 7.9, Theorem 7.10, and the finiteness/normalization results in
Sections 5--6. The restrictions on shared derivative pairs \((5,10)\),
\((10,15)\) imply more sparsity than a title-only comparison reveals.
Together with CLO they defeat novelty of the campaign's initial
five-term bound. The reported three-recycled-root computation is not a
monomial-count theorem; no complete public replay of its reported
\(3^{17}\) assignments was retrieved. The current arXiv metadata
supplies no journal reference.

\textbf{{[}Marashdeh{]} Mohammad F. Marashdeh, \emph{A descent-set
obstruction for the Casas-Alvero conjecture}.}
\href{https://arxiv.org/abs/2608.14726v1}{arXiv:2608.14726v1},
\href{https://arxiv.org/html/2608.14726v1}{full text}, 12 August 2026.
The record exposes only v1. Relevant comparisons concern support
stratification, triangular elimination of coefficients, descent-count
determinants, and the two-element-support argument. These are not a
degree-20 seven-total-term theorem. The campaign does not use a reported
equation count as a proof that the remaining incidence system is empty.
The preprint and its statements are distinct from external refereeing or
a priority clearance.

\textbf{{[}Bothmer-et-al{]} Hans-Christian Graf von Bothmer, Oliver
Labs, Josef Schicho, Christiaan van de Woestijne, \emph{The Casas-Alvero
conjecture for infinitely many degrees}.}
\href{https://arxiv.org/abs/math/0605090v2}{arXiv:math/0605090v2}, 25
June 2007; first submission 3 May 2006. \emph{Journal of Algebra} 316
(2007), 224--230, DOI
\href{https://doi.org/10.1016/j.jalgebra.2007.06.017}{10.1016/j.jalgebra.2007.06.017}.
The weighted projective formulation and good-prime transfer, including
Propositions 2.1--2.2 and 2.6, explain why geometric special-fibre
emptiness can prove a characteristic-zero result and propagate degree
\(n\) to \(np^e\). Prime-field rational-point absence or an empty
unsaturated affine chart does not meet those hypotheses.

\textbf{{[}Draisma-deJong{]} Jan Draisma, Johan P. de Jong, \emph{On the
Casas--Alvero conjecture}.} \emph{EMS Newsletter} 80 (June 2011),
29--33;
\href{https://ems.press/content/serial-issue-files/13698}{original issue
PDF}. \textbf{Read with the authors'
\href{https://math-unibe.ch/jdraisma/publications/erratumcasasalvero.pdf}{August
2011 erratum}.} Theorem 7's original \(4p^e\) statement omitted the bad
prime 5. The erratum explicitly withdraws the resulting degree-20
conclusion: the characteristic-5 quartic resultant equations have the
\(b=0\) solutions needed to defeat that argument. This is a correction
of a reduction argument, not a characteristic-zero counterexample.

\textbf{{[}Gasull{]} Armengol Gasull, \emph{A Primer on Resultants and
Their Applications}.}
\href{https://link.springer.com/article/10.1007/s44425-026-00047-6}{Publisher
full text}, DOI 10.1007/s44425-026-00047-6, published 28 May 2026.
Section 3.4's smallest-open-degree statement names 24 and cites
{[}Draisma-deJong{]}. The article's own Proposition 7 proves degrees 4
and 5. The cited erratum defeats the degree-20 consequence of the
original reference, and no replacement degree-20 proof appears in this
citation chain. That mathematical scope was checked directly.
``Inherited uncorrected statement'' is an inference about the
explanation, not a verified account of authorial intent.

\textbf{{[}Schaub-Spivakovsky-badprimes{]} Daniel Schaub, Mark
Spivakovsky, \emph{A description of and an upper bound on the set of bad
primes in the study of the Casas-Alvero Conjecture}.}
\href{https://arxiv.org/abs/2411.13967v1}{arXiv:2411.13967v1},
\href{https://arxiv.org/html/2411.13967v1}{full text}, 21 November 2024.
The full-scope regular-sequence/Macaulay-matrix formulation is relevant
to the unrestricted goal. Its upper bound on bad primes assumes the
characteristic-zero conjecture in the degree concerned; this assumption
cannot be used circularly to prove that degree.

\textbf{{[}Schaub-Spivakovsky-note{]} Daniel Schaub, Mark Spivakovsky,
\emph{A note on the Casas-Alvero Conjecture}.}
\href{https://arxiv.org/abs/2312.08742v7}{arXiv:2312.08742v7},
\href{https://arxiv.org/html/2312.08742v7}{full text}, 11 February 2025;
first submission 14 December 2023. Theorem 5's nonredundancy result for
high derivative resultants is weaker than a regular-sequence assertion.
Its reference to Ghosh is not an independent substitute for the missing
proof steps. Cite this version when discussing its added remarks.

\textbf{{[}Berger{]} Laurent Berger, \emph{The Weierstrass preparation
theorem and resultants of p-adic power series}.}
\href{https://arxiv.org/abs/1910.05319v2}{arXiv:1910.05319v2}, 4
November 2019; first submission 11 October 2019.
\href{https://perso.ens-lyon.fr/laurent.berger/articles/article33.pdf}{Author-hosted
PDF}. Section 1, including Corollary 1.2, supplies a precise
complete-ring preparation statement. In the campaign's row-4 local
argument, a power series reducing coefficientwise to \(z^2\) over a
complete unramified DVR is a unit times a monic distinguished quadratic.
The campaign separately proves the hypotheses, convergence, and coverage
of actual ramified points.

\textbf{{[}Lu-claim{]} Zhipeng Lu, \emph{Casas-Alvero conjecture in
computational algebraic geometry}.}
\href{https://arxiv.org/abs/1707.04754v6}{arXiv:1707.04754v6}, 16 March
2021; first submission 15 July 2017. The current abstract claims a full
proof via regular sequences. It was not independently certified in this
campaign. It is enough to disprove the bibliographic assertion that
Ghosh is the only author with a claimed full proof; it is not used as
evidence that the conjecture is settled.

\textbf{{[}Pereira{]} Rajesh Pereira, \emph{Matrix-theoretical
derivations of some results of Borcea--Shapiro on hyperbolic
polynomials}.} \emph{C. R. Acad. Sci. Paris, Ser. I} 341 (2005),
651--653.
\href{https://www.numdam.org/item/CRMATH_2005__341_11_651_0.pdf}{Primary
PDF}, DOI
\href{https://doi.org/10.1016/j.crma.2005.10.002}{10.1016/j.crma.2005.10.002}.
Accepted after revision 27 September 2005; available online 7 November
2005. Definitions 3--4 and Proposition 5, on p.~652, are exactly the
material used in the all-degree section: every complex matrix has a unit
trace vector, and compression to its orthogonal complement has
characteristic polynomial \(p_A'/n\) for the convention
\(p_A(X)=\det(XI-A)\). This statement includes nonnormal matrices,
despite the article's hyperbolic-polynomial title. The 2005 article
explicitly recalls the result from the author's earlier
\emph{Differentiators and the geometry of polynomials}, \emph{J. Math.
Anal. Appl.} 285 (2003), 336--348, Theorem 2.5. The 2005 title should
not be replaced with that distinct 2003 title. Primary PDF and its usage
were checked directly on 23 September 2026; the HTML landing page
failed, but the PDF was accessible.

\subsection{Provenance and unresolved source
gaps}\label{provenance-and-unresolved-source-gaps}

\begin{enumerate}
\def\labelenumi{\arabic{enumi}.}
\item
  \textbf{ProofAtlas overlap.} The
  \href{https://www.proofatlas.ai/collaboration/casas-alvero-conjecture/}{project's
  own Casas--Alvero page} was inspected in the campaign's bounded
  priority rounds. It reports degree-20 certificate/Hensel work and
  systems called A1/A2/A3, but the retrieved material did not provide
  enough exponents, full integer systems, or certificates to compare
  with this campaign's C family. Those labels are not interchangeable
  with ours. This is primary evidence of reported activity, not a
  verified theorem. No priority clearance, outreach, or assumption of
  non-overlap follows.
\item
  \textbf{Unavailable thesis.} The primary
  \href{https://etd.lib.nycu.edu.tw/cgi-bin/gs32/hugsweb.cgi?o=dnthucdr&s=id\%3D\%22G021090215100\%22.&searchmode=basic}{NTHU
  catalog record for Shih Cheng Pang's 2022 master's thesis, \emph{On
  the Casas-Alvero Conjecture}}, was located, but the full text was not
  obtained. No overlap conclusion is made.
\item
  \textbf{Unverified criticism lead.} A search lead involving the
  prequel's Lemma 5.4 was not recovered as a sufficiently complete
  primary mathematical argument. This review does not attribute that
  criticism to an identified author, repeat it as an established result,
  or use it to discredit the accepted paper. The explicit audit evidence
  stands on its own calculations.
\item
  \textbf{Chellali and small bad primes.} The earlier literature
  inventory consulted
  \href{https://arxiv.org/abs/1211.2059}{arXiv:1211.2059} and
  primary-manuscript material concerning degree \(5p^e\) and exceptional
  primes. A fresh bibliographic re-verification was not needed for this
  review's proof: the degree-5 characteristic-2 seed is displayed and
  checked directly, and the degree-4 characteristic-5 obstruction has
  the primary authors' explicit erratum. Do not promote this pointer to
  a newly audited complete theorem or fill missing metadata by
  guesswork.
\item
  \textbf{Exact certificate versus priority.} The initial audit package
  and its \href{evidence/AUDIT_NOVELTY_CORRECTION.md}{unchanged
  correction} are both preserved. The correction retracts candidate
  novelty of the five-total-term lower bound while preserving its valid
  certificate. Its historical ``no six-term theorem'' statement and the
  predecessor's unresolved-C status are superseded by
  \href{evidence/seven-terms/PROOF.md}{the seven-term proof}. Earlier
  arithmetic is not silently rewritten.
\item
  \textbf{What was refreshed.} This final pass reread the Ghosh record,
  relevant local source definitions, the Draisma--de Jong erratum, the
  Gasull citation chain, arXiv version records, the publisher's
  Ghosh-finiteness acceptance entry, and the thesis metadata. It did not
  start an exhaustive new literature survey. Earlier exact-family
  queries and their limitations remain in the bundled sparse prior-art
  notes. No new degree-20 proof was identified in the bounded checked
  material; the stronger universal claim that no such proof or comment
  exists is not made.
\item
  \textbf{Proof foundations.} Elementary valuation comparisons,
  unit-Jacobian bounds, simple-root uniqueness, finite-dimensional
  local-algebra arguments, and the explicit polynomial identities in
  this review are proved or reduced to replayable identities in the
  corresponding evidence. Their use does not constitute a claim of new
  general theory. The review does not rely on external publication or a
  formal proof assistant to validate those arguments.
\end{enumerate}

\end{document}
